\documentclass[11pt]{article}
\usepackage{microtype,amsmath,amssymb,amsthm,mathtools,booktabs,array}
\usepackage[margin=1.05in]{geometry}
\usepackage[colorlinks=true,linkcolor=blue!45!black,citecolor=blue!45!black,urlcolor=blue!45!black]{hyperref}
\usepackage{xcolor}
\hypersetup{pdftitle={The dimension-degree boundary for Gaussian moment vanishing},pdfsubject={Working mathematical manuscript}}
\newtheorem{theorem}{Theorem}[section]
\newtheorem{proposition}[theorem]{Proposition}
\newtheorem{lemma}[theorem]{Lemma}
\newtheorem{corollary}[theorem]{Corollary}
\theoremstyle{definition}\newtheorem{definition}[theorem]{Definition}
\theoremstyle{remark}\newtheorem{remark}[theorem]{Remark}
\newcommand{\E}{\mathbb E}
\newcommand{\C}{\mathbb C}
\newcommand{\tr}{\operatorname{tr}}
\newcommand{\wt}{\operatorname{wt}}
\newcommand{\G}{\mathsf G}
\DeclareMathOperator{\Span}{span}
\title{The dimension--degree boundary\\for Gaussian moment vanishing}
\author{}
\date{Working draft, 17 November 2026}
\begin{document}
\maketitle
\begin{center}\begin{minipage}{0.91\textwidth}\small
\textbf{Draft status.} This revision supplies explicit algebraic and formal
lemmas, a complete proof of the mixed cubic chart, and exhaustive routing
tables. A general formal leading-coefficient criterion now closes the
rank-one quadratic branch, including its escaping and degenerate cases.
A resolved-curve proof now supplies the full logarithmic estimate,
including two irregular ends. The affine-linear-coordinate classification
is complete within this draft. A common gradient argument also settles
smooth leading forms in every degree and the coercive ordinary-node
branches. The finite-critical no-frame members of all three ordinary-node
families are now excluded by one action identity and an explicit
regular-block rank argument. Critical conics and hyperbolas are also settled by one all-degree
formal-frame theorem. Center axes and the other geometric
no-frame exclusions still need
complete proofs. The
dimension--degree conclusion is stated using the retained ternary theorem;
this manuscript is not yet its submission-ready standalone proof.
\end{minipage}\end{center}

\begin{abstract}
The Gaussian Moments Conjecture asks whether the vanishing of every
positive moment of a complex polynomial persists after multiplication by
an arbitrary fixed polynomial. The answer depends on both the number of
variables and the degree. Combining the ternary-cubic theorem with
Wilson's two-variable theorem and Long's counterexamples gives the exact
boundary: the conjecture holds universally precisely when there are at
most two variables, the degree is at most two, or there are three variables
and the degree is at most three. In the last case the explicit cutoff is
$m>3\deg Q$, and the coefficient three is optimal. We explain why these
statements fit together, give short proofs of the elementary parts, and
isolate the cubic reductions that carry the new argument.
\end{abstract}

\section{What vanishing moments tell us}

Let $X_1,\ldots,X_n$ be independent standard real Gaussian random
variables. We allow a polynomial $P(X)$ to have complex coefficients.
Suppose that
\begin{equation}\label{eq:pure}
  \E[P(X)^m]=0\qquad\text{for every integer }m\geq1.
\end{equation}
The Gaussian Moments Conjecture asks whether this forces
\begin{equation}\label{eq:safe}
  \E[Q(X)P(X)^m]=0\qquad\text{for all sufficiently large }m
\end{equation}
for every fixed polynomial $Q$. The threshold may depend on both $P$
and $Q$.

The complex coefficients matter. For real $P$, the second moment in
\eqref{eq:pure} forces $P=0$. For complex $P$, cancellation can make
every positive moment vanish even when $P$ is far from zero. For example,
$(X_1+iX_2)/\sqrt2$ has this property. The question is whether a fixed
polynomial multiplier can keep exposing a nonzero moment at arbitrarily
large powers.

We call \eqref{eq:pure} \emph{purity} and \eqref{eq:safe}, for every
$Q$, \emph{safety}. These names distinguish the hypothesis from the
conclusion; safety without purity is a separate question. In the original
formulation of Derksen, van den Essen and Zhao~\cite{DVEZ}, the assertion
that every pure polynomial is safe was a conjecture in every dimension.
Long's examples show why the dimension cannot be ignored~\cite{Long}.

Here the question is more precise. For integers $n\geq1$ and $d\geq0$,
let $\G(n,d)$ mean that every pure complex polynomial of total degree at
most $d$ in $n$ standard Gaussian variables is safe.

\begin{theorem}[Dimension--degree boundary]\label{thm:boundary}
Using the ternary-cubic theorem stated in Theorem~\ref{thm:cubic},
\[
 \boxed{\quad \G(n,d)\ \Longleftrightarrow\
 n\leq2\ \text{or}\ d\leq2\ \text{or}\ (n,d)=(3,3).\quad}
\]
\end{theorem}

This statement classifies the ranges in which purity \emph{always}
implies safety. It does not classify the individual safe polynomials
outside those ranges. A counterexample in four variables does not say
that all four-variable polynomials behave badly.

\begin{center}
\begin{tabular}{@{}lccc@{}}
\toprule
 & $d\leq2$ & $d=3$ & $d\geq4$\\
\midrule
$n=1,2$ & always safe & always safe & always safe\\
$n=3$ & always safe & always safe & failures exist\\
$n\geq4$ & always safe & failures exist & failures exist\\
\bottomrule
\end{tabular}
\end{center}

The only positive case needed beyond the existing binary and quadratic
theorems is a cubic in three variables. Its quantitative form is useful:

\begin{theorem}[Ternary cubics; internal theorem]\label{thm:cubic}
Let $P\in\C[X_1,X_2,X_3]$ have total degree at most three and satisfy
\eqref{eq:pure}. Then, for every polynomial $Q$,
\begin{equation}\label{eq:cutoff}
 \E[QP^m]=0\qquad\text{whenever }m>3\deg Q.
\end{equation}
No smaller constant can replace three uniformly over this class.
\end{theorem}

The sharpness is visible in a three-term example, not in a large
calculation. We prove it in Section~\ref{sec:shear}. The difficult part
is showing that an arbitrary pure ternary cubic can be reduced to the
families where a direct argument works. Section~\ref{sec:cubic} separates
that problem from the boundary deduction.

\subsection*{Attribution and the scope of the paper}
Wilson's Theorem~1 supplies the unrestricted two-variable
result~\cite{Wilson}. Liu's Theorems~2.1 and~2.2 give the quadratic classification and optimal
mixed cutoff~\cite{Liu}.
Long's Proposition~4.1 and Theorem~5.1 supply the cubic four-variable
and quartic three-variable counterexamples, respectively~\cite{Long}. These are
credited inputs, not new examples or new binary and quadratic theorems.
The intended new contribution is the ternary-cubic theorem, its readable
proof, and the resulting exact dimension--degree classification.
Wilson's and Long's cited works are public preprints; this draft does
not describe them as peer-reviewed journal articles.

Wilson already obtains the classification by dimension without a degree
restriction. The proposed contribution here is the finer dimension--degree
boundary through the ternary-cubic result, not a new claim to that
dimension-only classification. A search of the selected primary records
located the credited binary, quadratic and counterexample results but
did not establish priority for the assembled boundary. No first-proof
or exhaustive novelty claim is made in this draft.


\section{A few Gaussian identities}

All our generating functions are formal power series. We never need to
assume that the integral of $e^{tP(X)}$ converges. Define
\begin{equation}\label{eq:phi}
 \Phi_Q(t)=\sum_{m\geq0}\E[QP^m]\frac{t^m}{m!}.
\end{equation}
Then $P$ is pure exactly when $\Phi_1=1$, and is safe exactly when
every $\Phi_Q$ is a polynomial. Thus a degree bound on $\Phi_Q$ gives
an explicit vanishing threshold.

For two independent real Gaussian variables, put
\[
 Z=\frac{X_1+iX_2}{\sqrt2},\qquad
 W=\frac{X_1-iX_2}{\sqrt2}.
\]
Although $Z,W$ are convenient polynomial coordinates, they are not
independent random variables: $W=\overline Z$. Their moments are
\begin{equation}\label{eq:wick}
 \E[Z^aW^b]=\delta_{ab}a!.
\end{equation}
Rotating $(X_1,X_2)$ proves the vanishing when $a\ne b$; when $a=b$,
$ZW=(X_1^2+X_2^2)/2$ has moments $a!$. An additional independent
standard Gaussian $T$ has $\E[T^{2j}]=(2j-1)!!$ and zero odd moments.

\begin{lemma}[One sign of weight]\label{lem:weight}
Give $Z,W,T$ weights $1,-1,0$, respectively. If every monomial of $P$
has strictly positive weight, or every monomial has strictly negative
weight, then $P$ is pure and
\[
 \E[QP^m]=0\qquad(m>\deg Q).
\]
\end{lemma}
\begin{proof}
In the positive case, every monomial of $P^m$ has weight at least $m$.
A monomial of $Q$ has weight at least $-\deg Q$. For $m>\deg Q$ their
product therefore has nonzero weight, and \eqref{eq:wick} kills it.
The negative case follows by exchanging $Z$ and $W$.
\end{proof}

Complex orthogonal changes of polynomial coordinates preserve the
Gaussian moment functional. This follows from its source identity
$\E[e^{u^TX}]=e^{u^Tu/2}$. We use this algebraic invariance; we do not
interpret a complex coordinate change as a real change of integration
contour.

\begin{lemma}[Linear circular contraction]\label{lem:contraction}
Let $a,b$ be constants and let $F(t,z)$ be a formal power series in $t$
whose coefficients are polynomials in $z$. Then
\begin{equation}\label{eq:contraction}
 \E\!\left[F(t,Z)e^{t(a+bZ)W}\right]
   =\frac{1}{1-bt}\,
       F\!\left(t,\frac{at}{1-bt}\right).
\end{equation}
\end{lemma}
\begin{proof}
It suffices, coefficient by coefficient in $t$, to take $F=Z^j$.
Using \eqref{eq:wick}, the left side is
\[
 \sum_{k\geq j}t^k[z^k]z^j(a+bz)^k
 =\sum_{r\geq0}\binom{j+r}{j}a^j b^r t^{j+r}
 =\frac{(at)^j}{(1-bt)^{j+1}}.
\]
Every fixed coefficient uses only finitely many terms, which also
justifies the statement for $t$-dependent $F$.
\end{proof}

\section{Quadratics in any number of variables}\label{sec:quadratics}

We include a short proof of the positive quadratic result. It explains
why degree two is safe in every dimension without introducing the
finite-moment thresholds that are unnecessary for the boundary.
The classification and sharper surrounding results are due to
Liu~\cite[Theorem~2.1]{Liu}.

\begin{proposition}\label{prop:quadratic}
Every pure complex quadratic in $n$ standard Gaussian variables is safe.
In fact
\[
 \E[QP^m]=0\qquad\left(m>\left\lfloor n/2\right\rfloor\deg Q\right).
\]
\end{proposition}
\begin{proof}
Write $P=X^TAX+\ell^TX+c$, with $A=A^T$, and put
$\Sigma(t)=(I-2tA)^{-1}$ and $D(t)=\det(I-2tA)$. Completing the
square formally gives
\begin{equation}\label{eq:quad}
 \Phi_1(t)=D(t)^{-1/2}
   \exp\!\left(ct+\frac{t^2}{2}\ell^T\Sigma(t)\ell\right).
\end{equation}
If $\Phi_1=1$, then $D=e^S$ for a rational function $S$ regular at
zero. Hence $D'/D=S'$. A rational derivative has zero residue at every
finite pole, whereas $D'/D$ has the positive multiplicity as residue at
every zero of $D$. Thus $D$ has no zero, so $D=1$. It follows that
$A$ is nilpotent and
\[
 c=0,\qquad \ell^TA^j\ell=0\quad(j\geq0).
\]
Conversely these identities make \eqref{eq:quad} equal to one.

With an additional source $u$, the same calculation now gives
\begin{equation}\label{eq:quadsource}
 \E[e^{tP+u^TX}]
   =\exp\!\left(u^T\mu(t)+\tfrac12u^T\Sigma(t)u\right),
 \qquad \mu(t)=t\Sigma(t)\ell.
\end{equation}
Both $\mu$ and $\Sigma$ are polynomials in $t$. Let $s$ be the dimension
of $\Span\{\ell,A\ell,A^2\ell,\ldots\}$, with $s=0$ if $\ell=0$.
This space is totally isotropic, since
$(A^i\ell)^TA^j\ell=\ell^TA^{i+j}\ell=0$. Nondegeneracy of the
ambient form gives $s\leq h:=\lfloor n/2\rfloor$. Nilpotence gives
$A^s\ell=0$, so $\deg_t\mu\leq s$. Also
$\deg_t\Sigma\leq n-1\leq2h$.

For a monomial of degree $q$ in $X$, a term in
\eqref{eq:quadsource} uses $p$ covariance factors and $q-2p$ mean
factors. Its $t$-degree is at most
$2hp+h(q-2p)=hq$. Source differentiation and linearity prove the result.
\end{proof}

\section{Two variables, with no degree restriction}

The needed input is Wilson's support theorem~\cite[Theorem~1]{Wilson}:
if a polynomial in two Gaussian variables is pure, then in the coordinates
$Z,W$ every monomial has weight $\deg_Z-\deg_W$ of one strict sign.
Lemma~\ref{lem:weight} immediately gives the threshold $m>\deg Q$.
This proves $\G(2,d)$ for every $d$.

The one-variable conclusion follows as well. Embed a one-variable
polynomial into two variables by ignoring the second. A multiplier in
the larger space can be averaged in that unused variable. In particular,
safety in two variables implies safety in one. No extension of Wilson's
theorem to three variables is being assumed here.

\section{The two failures that determine the boundary}\label{sec:failures}

Long's examples~\cite{Long} have the same marked generating function.
Writing them side by side makes both the threshold and the failure of
eventual vanishing transparent.

\subsection{A quartic in three variables}
Use $Z,W,T$ as above and set
\[
 P_3=(1+Z)\left(W-\frac{2+Z}{2}T^2\right),\qquad Q_3=Z.
\]
The total degree is four. For any polynomial $A$, first average over
$T$, then apply Lemma~\ref{lem:contraction} with $a=b=1$. Put
$z_*=t/(1-t)$. The result is
\begin{align*}
 \E[A(Z)e^{tP_3}]
 &=\frac{A(z_*)}{1-t}
       \left(1+t(1+z_*)(2+z_*)\right)^{-1/2}\\
 &=A\!\left(\frac{t}{1-t}\right),
\end{align*}
because the expression inside the square root is $(1-t)^{-2}$.
The root with constant term one cancels the factor $(1-t)^{-1}$.

\subsection{A cubic in four variables}
Use two independent circular pairs $(Z_1,W_1)$ and $(Z_2,W_2)$ and set
\[
 P_4=(1+Z_2)\bigl(W_1(1-Z_1)+W_2\bigr),\qquad Q_4=Z_2.
\]
The total degree is three. For a formal parameter $s$, circular
contraction gives
\[
 \E[e^{sW_1(1-Z_1)}]=\frac1{1+s}.
\]
Substitute $s=t(1+Z_2)$, then contract the second pair. Again
\[
 \E[A(Z_2)e^{tP_4}]
  =\frac{A(z_*)}{1-t}\frac{1}{1+t(1+z_*)}
  =A\!\left(\frac{t}{1-t}\right).
\]

In either example, $A=1$ gives purity, while $A=z$ gives
\[
 \Phi_Q(t)=\frac{t}{1-t},\qquad
 \E[QP_j^m]=m!\ne0\quad(m\geq1).
\]
Thus one linear multiplier sees a nonzero moment at every positive
power. These examples disprove $\G(3,4)$ and $\G(4,3)$.

\section{Why the cubic cutoff is three}\label{sec:shear}

The factor three in Theorem~\ref{thm:cubic} has a concrete source.
Consider
\begin{equation}\label{eq:shearfamily}
 P=aW+(c+dZ)T-\frac{Z(c+dZ)^2}{2a},\qquad a\ne0.
\end{equation}
For sources $\alpha,\beta,\gamma$, first contract $W$ and then
average over $T$. If $u(z)=c+dz$, the exact identity is
\begin{equation}\label{eq:shearsource}
\begin{split}
 \E[e^{tP+\alpha Z+\beta W+\gamma T}]
  =\exp\biggl(&\alpha(at+\beta)+\frac{\gamma^2}{2}
       +\gamma t\,u(at+\beta)\\[-2pt]
       &-\frac{\beta t}{2a}u(at+\beta)^2\biggr).
\end{split}
\end{equation}
Setting the sources to zero proves purity. Give the sources
$\alpha,\beta,\gamma$ weights $1,3,2$, respectively. Every term in
the exponent has $t$-degree at most its source weight. The same is true
after exponentiation. It follows that
\[
 \deg_t\Phi_{Z^iW^jT^k}\leq i+3j+2k\leq3(i+j+k).
\]
Thus every member of \eqref{eq:shearfamily} satisfies
\eqref{eq:cutoff}.

Now take the simplest member,
\[
 P_*=W+ZT-\frac12 Z^3.
\]
With only the $W$ source left, \eqref{eq:shearsource} becomes
\[
 \E[e^{tP_*+\beta W}]
   =\exp\!\left(-\frac{\beta t(t+\beta)^2}{2}\right).
\]
The highest $t$-degree in its $\beta^D$ coefficient comes only from
$(-\beta t^3/2)^D/D!$. Therefore, for every $D\geq1$,
\begin{equation}\label{eq:sharp}
 \E[W^D P_*^{3D}]=(3D)!\left(-\frac12\right)^D\ne0.
\end{equation}
Since $\deg W^D=D$, any proposed uniform coefficient smaller than
three would already fail on this family. This establishes the sharpness
part of Theorem~\ref{thm:cubic} independently of the general reduction.

\section{Reducing an arbitrary ternary cubic}\label{sec:cubic}

The preceding sections prove the easy parts of the boundary and explain
the sharp constant. The remaining issue is a cubic with no preferred
coordinate. Write it in homogeneous pieces as
\[
 P=F_3+F_2+F_1+F_0.
\]
We want a direction in which $P$ is only affine linear. The direction
must also be isotropic for the Gaussian bilinear form; an arbitrary
linear change would lose the moment identities.

\begin{lemma}[The algebraic frame test]\label{lem:frame}
There are Gaussian coordinates $(Z,W,T)$ in which $P$ is affine linear
in $W$ if and only if there is a nonzero vector $v\in\C^3$ such that
\begin{equation}\label{eq:frame}
 v^Tv=0,\qquad \nabla F_3(v)=0,\qquad F_2(v)=0.
\end{equation}
\end{lemma}
\begin{proof}
Polarization gives
\[
 D_v^2P(x)=2\nabla F_3(v)^Tx+2F_2(v).
\]
The last two conditions in \eqref{eq:frame} say exactly that this
polynomial vanishes. Complete the nonzero null vector $v$ to a Witt
basis: choose a null partner with pairing one, and a unit vector in
their orthogonal complement. In these coordinates $D_v$ is the $W$
derivative. Hence $D_v^2P=0$ says that no second or higher power of
$W$ occurs. The converse reads off $v$ from the $W$ direction.
\end{proof}

The hard reduction is an existence statement, not this algebraic test.

\begin{proposition}[Frame extraction; consolidation step A]\label{prop:frameextract}
If a ternary cubic is pure, then a vector satisfying
\eqref{eq:frame} exists.
\end{proposition}

The retained proof divides according to the projective cubic $F_3=0$.
Over $\C$, the list is short: a smooth cubic, an irreducible nodal or
cuspidal cubic, a line with a smooth conic, three distinct lines, or a
repeated line. The first group has no leading singular direction at
all. In the singular cases the issue is whether the singular direction
is null and whether $F_2$ vanishes there. Cases failing this test are
excluded by their period behavior. Once a direction passes the test,
Lemma~\ref{lem:frame} supplies the desired coordinates.

This paragraph is a roadmap, not a replacement for the exclusion proof.
The standalone version must show that the period argument accounts for
degenerate critical points, escaping critical points and lower-degree
terms. The exact cases and source notes are recorded in the companion
proof audit. No numerical sample or finite moment computation proves
Proposition~\ref{prop:frameextract}.

\begin{proposition}[Affine-linear coordinate; consolidation step B]\label{prop:wlinear}
Let
\[
 P=W(a+bZ+cT+dZ^2+eZT+fT^2)+G(Z,T),\qquad\deg G\leq3.
\]
If $P$ is pure, it is safe and satisfies $m>3\deg Q$.
\end{proposition}

The proof below reduces the pure cases to the one-sign argument,
the quadratic calculation, and the shear mechanism of
Section~\ref{sec:shear}. Some excluded cases first reduce to a
one-variable polynomial exponential period. The proof must retain
the determinant factor from eliminating a Gaussian variable; omitting
it can turn a decaying period into an apparent constant. It must also
use the actual period selected by the Gaussian moment functional,
rather than every solution of an enlarged differential equation.

One substantial part of this classification has a short algebraic proof.
It is useful both as a proof simplification and as a check on what the
remaining period analysis must establish.

\subsection{A complete algebraic subcase}\label{sec:algebraic}
Suppose
\[
 P=W h(Z)+v(Z)T^2+u(Z)T+r(Z),
\]
where the coefficient degrees make $P$ cubic. Thus
$\deg h\leq2$, $\deg v\leq1$, $\deg u\leq2$, $\deg r\leq3$.
If $g=th(g)$ and $J=1-th'(g)$, formal circular contraction followed
by the quadratic Gaussian identity gives
\begin{equation}\label{eq:subcase}
 \Phi_1=
 \frac{\exp\!\left(tr(g)+t^2u(g)^2/(2K)\right)}{J\sqrt K},
 \qquad K=1-2tv(g).
\end{equation}
The circular identity here is the one-variable formal Lagrange formula;
it also follows directly from coefficient extraction. For completeness,
for any numerator $A$, its formula is
\[
 \E[A(Z)e^{tWh(Z)}]=\frac{A(g)}{1-th'(g)}.
\]
Indeed the left side is $\sum_{m\geq0}t^m[z^m]A(z)h(z)^m$;
the same expression is the formal residue of $A(z)/(z-th(z))$,
whose unique small pole is $g$. The identity remains valid for the
$t$-dependent numerator in \eqref{eq:subcase}.

If $h(0)\ne0$, change variables formally from $t$ to $z=g(t)$ and put
\[
 A=h(z),\quad B=h(z)-2zv(z),\quad C=h(z)-zh'(z).
\]
The equation $\Phi_1=1$, after squaring, equates the exponential of a
rational function with a rational function. The residue argument used
in Section~\ref{sec:quadratics} forces
\begin{equation}\label{eq:subcaseidentities}
 C^2B=A^3,\qquad 2Br+zu^2=0.
\end{equation}
These identities first force $\deg h\leq1$. To see this, if the distinct
nonzero roots of $h$ are $\alpha_i$ with multiplicities $m_i$, then
\[
 \frac Ch=1-z\sum_i\frac{m_i}{z-\alpha_i}.
\]
When $\deg h\geq2$, the reduced numerator has the same positive degree
as its denominator, has leading coefficient $1-\deg h$, and is nonzero
at each $\alpha_i$. It therefore has a root away from the roots of $h$,
contradicting $C^2\mid h^3$.

If $h=a\ne0$, \eqref{eq:subcaseidentities} gives $v=0$ and
$r=-zu^2/(2a)$. The cubic degree bound forces $\deg u\leq1$, giving
exactly \eqref{eq:shearfamily}. If $h=a+bz$ with $ab\ne0$, it gives
\[
 v=\frac{h-h^3/a^2}{2z}.
\]
This has a nonzero quadratic coefficient, contradicting $\deg v\leq1$.
Thus this second branch cannot be a cubic.

If $h(0)=0$, the fixed point is $g=0$. Writing
$b=h'(0)$, $v_0=v(0)$, $u_0=u(0)$ and $r_0=r(0)$, purity in
\eqref{eq:subcase} forces
\[
 (1-bt)^2(1-2v_0t)=1,\qquad
 tr_0+\frac{t^2u_0^2}{2(1-2v_0t)}=0.
\]
Hence $b=v_0=u_0=r_0=0$. Every surviving monomial has strictly positive
$Z$-minus-$W$ weight, so Lemma~\ref{lem:weight} applies.

This proves Proposition~\ref{prop:wlinear} on the displayed subcase
without a case table, an analytic limit, or a finite-moment bound.
The next subsection removes $T^3$ algebraically and also treats the
linear $T$ coefficient of $W$. A general formal valuation argument then
closes both the mixed $WT/WZ^2$ chart and the rank-one $WT^2$ family. The subsequent case table is exhaustive. The logarithmic estimate
is proved on a resolved curve, retaining both its comparison theorem
and its quantitative contour bounds.

\subsection{Removing the cubic height without asymptotics}\label{sec:airyformal}

The extra term $T^3$ can be dealt with algebraically. This avoids
using an asymptotic formula for a cubic Gaussian integral.

\begin{lemma}[An algebraic Riccati obstruction]\label{lem:airyriccati}
The equation $f_x=x-f^2$ has no solution algebraic over $\C(x)$.
\end{lemma}
\begin{proof}
At a finite point, a finite branch of $f$ satisfies a regular analytic
first-order equation and cannot branch. At a pole, put $g=1/f$.
Then $g_x=1-xg^2$, again a regular equation; a zero of $g$ is simple.
Thus the normalized algebraic curve of $\C(x,f)$ is unramified over
every finite $x$. If its degree is $d$ and there are $r$ points over
infinity, Riemann--Hurwitz gives $2g-2=-d-r$. This forces $d=r=1$.
So $f$ would be rational. An integer-order pole of $f$ at infinity
makes $f^2$ dominate both $x$ and $f_x$; if $f$ is bounded there,
the term $x$ cannot cancel. Both possibilities are impossible.
\end{proof}

\begin{lemma}[Cubic Gaussian obstruction]\label{lem:cubicgaussian}
Let $A,B,C\in\C(z)$ be regular at zero, with $A(0)=C(0)=0$,
$B(0)=-1/2$ and $A\ne0$. Define the formal unit
\[
 I=\E_T\exp\bigl(AT^3+(B+1/2)T^2+CT\bigr).
\]
Then $I'/I$ is not rational in $z$.
\end{lemma}
\begin{proof}
Write $M_j=\E[T^j\exp(AT^3+(B+1/2)T^2+CT)]$ and $J=M_1$.
Gaussian integration by parts gives, for $j\geq0$,
\begin{equation}\label{eq:momentrec}
 3A M_{j+2}+2B M_{j+1}+C M_j+jM_{j-1}=0,
\end{equation}
where the last term is zero for $j=0$. Also
$I'=A'M_3+B'M_2+C'J$ and $J'=A'M_4+B'M_3+C'M_2$.
For clarity, the first two eliminations are
\[
 M_2=-\frac{2BJ+CI}{3A},\qquad
 M_3=\frac{(4B^2-3AC)J+(2BC-3A)I}{9A^2}.
\]
The case $j=2$ of \eqref{eq:momentrec} supplies $M_4$.

In an algebraic extension choose $a^3=3A$ and put
\[
 X=\frac{B^2-3AC}{a^4},\qquad F=\frac JI,
 \qquad f=-aF-\frac B{a^2}.
\]
Substitution in these three recurrences, using $a'=A'/a^2$, gives
\begin{align}
 \frac{I'}I&=U-aX'F,
 &U&=\frac{A'(2BC-3A)}{9A^2}-\frac{B'C}{3A},\label{eq:airyfirst}\\
 f'&=X'(X-f^2).\label{eq:airypullback}
\end{align}
These identities follow by differentiation and polynomial substitution;
they do not involve an integration contour or a limit of a Gaussian
integral. They can also be checked by substituting $A=a^3/3$ in
\eqref{eq:momentrec} before differentiating.

If $A$ has order $q>0$, then $X$ has a pole of order $4q/3$,
because $B^2-3AC$ is a unit. Thus $X'\ne0$ and $aX'$ is a nonzero
rational function. If $I'/I$ were rational, \eqref{eq:airyfirst}
would make $F$ rational and $f$ algebraic over $\C(z)$. The function
field $\C(z,a)$ is finite over $\C(X)$, so $f$ would be algebraic
over $\C(X)$. Equation~\eqref{eq:airypullback} contradicts
Lemma~\ref{lem:airyriccati}.
\end{proof}

\begin{proposition}[The thirteen-coefficient case]\label{prop:height}
Suppose a cubic has the form
\[
 P=Wh(Z)+sT^3+v(Z)T^2+u(Z)T+r(Z),
\]
where $s$ is constant. If it is pure, then $s=0$; it is one of the
pure cases described in Section~\ref{sec:algebraic}, and obeys
$m>3\deg Q$.
\end{proposition}
\begin{proof}
If $h(0)\ne0$, use $z=g$, $t=z/h(z)$ in the formal Lagrange
identity. The remaining cubic Gaussian series has
\[
 A=\frac{zs}{h},\quad B=-\frac12+\frac{zv}{h},\quad C=\frac{zu}{h}.
\]
Purity would give $I=H e^{-R}$, where
$H=(h-zh')/h$ and $R=zr/h$. If $s\ne0$, this has rational
logarithmic derivative, contradicting Lemma~\ref{lem:cubicgaussian}.
If $h(0)=0$, the fixed point is zero. Use instead
$A=ts$, $B=-1/2+tv(0)$, $C=tu(0)$, $H=1-th'(0)$ and $R=tr(0)$.
The same contradiction applies. Thus $s=0$ in both cases, and the
already proved algebraic subcase finishes the argument.
\end{proof}

\begin{corollary}[The affine coefficient of $W$]\label{cor:linearwt}
If $P=W(a+bZ+cT)+G(Z,T)$ is a pure cubic, it satisfies
$m>3\deg Q$. If $c\ne0$, it is in fact quadratic.
\end{corollary}
\begin{proof}
Put $q=1-bt$ and $g=at/q+(ct/q)T$. Circular contraction gives
\[
 \Phi_1=q^{-1}\E_T[e^{tG(g,T)}].
\]
The coefficient of $T^3$ in the exponent is
$A(t)=tG_3(ct/q,1)$. If it is nonzero, the cubic Gaussian
obstruction applies, whereas purity makes the remaining Gaussian
series equal to $q$, after removing a rational scalar exponential.
For $c\ne0$, the rational function $ct/q$ is nonconstant, so
$A=0$ forces $G_3=0$. Proposition~\ref{prop:quadratic} applies.
If $c=0$, use Proposition~\ref{prop:height}.
\end{proof}

\subsection{What purity gives in a differential module}\label{sec:seedbridge}

The following lemma is the precise passage from the formal Gaussian
functional to an analytic period. Its hypotheses must be verified
after each change of variables or localization.

\begin{lemma}[The seed functional]\label{lem:seedfunctional}
Let $M$ be a finite-dimensional differential module over $K=\C(t)$,
let $s\in M$, and suppose a $K$-linear horizontal map
$J:M\longrightarrow\C((t))$ satisfies $J(s)=1$. Then the differential
submodule $U=K\langle\partial_t\rangle s$ has a horizontal quotient
$L:U\longrightarrow K$ with $L(s)=1$.

At an ordinary base point, $L$ extends locally to an analytic horizontal
functional on $M$. For the period conclusion, assume explicitly that
$M$ is the generic exponential de Rham system of regular functions
$f,g$ on a smooth affine complex variety, where $f$ is the base
projection and $f,g$ are algebraically independent. Remove the finite
exceptional base set in the Hien--Roucairol comparison theorem.
Then the extension is represented by a flat rapid-decay cycle, and its
period evaluated on $s$ is identically one.
\end{lemma}
\begin{proof}
On the generators $s,\nabla s,\nabla^2s,\ldots$, the values of $J$
are $1,0,0,\ldots$. Therefore its restriction takes values in $K$.
Every rational relation between these generators is respected because
$J$ was defined on $M$, so this restriction is well defined and
horizontal. Its value at $s$ makes it surjective.

Choose a point where $M$, $U$ and $L$ are ordinary, extend the fiber
functional from $U$ to $M$, and solve the dual analytic ODE.
Uniqueness preserves its restriction to $U$. The perfect period
pairing and solution representation in
Hien--Roucairol~\cite[Theorems~2.4 and~3.7]{HienRoucairol} give the
asserted cycle. A scalar rational normalization is an invertible gauge
away from its zeros and poles and must be retained in the period.
\end{proof}

There are three distinct verifications here. First, integration by
parts must make $J$ descend to the exact displayed de Rham quotient.
Second, every inverted spatial coordinate must be a valid formal
Laurent unit on the selected branch. Third, the comparison theorem
requires a smooth affine total space and regular functions after
removing the finite set of poles in the parameter. The theorem gives
a representation, not a bound at infinity. In particular, it gives
no shrinking estimate on its own.

There is also a useful purely formal way to obtain a contradiction.

\begin{lemma}[A full formal frame and a positive seed]\label{lem:positiveformal}
Assume the hypotheses of the first paragraph of
Lemma~\ref{lem:seedfunctional}. Suppose that, after a finite ramification
$r=t^{-1/N}$ at infinity, $M$ has a full formal frame of rank-one
exponential factors. For every factor with zero exponential part,
assume its horizontal functional takes $s$ into $r\C[[r]]$.
Then $J(s)=1$ is impossible.
\end{lemma}
\begin{proof}
Enlarge the ramification to include the rational powers in the frame.
The module over $\C((r))$ is then a direct sum of rank-one modules,
so the quotient on $U$ extends to a horizontal map from the whole
module to $\C((r))$. A nonzero exponential factor has no such map:
the logarithmic derivative of a Laurent series has at most a simple
logarithmic pole and cannot equal the derivative of a nonzero
negative-power exponential part. On the zero part, all horizontal
functionals are constant linear combinations of the supplied frame.
Their values on $s$ have strictly positive $r$-order and cannot sum
to one.
\end{proof}

The word \emph{full} is essential. Counting some critical points or
writing a few scalar asymptotic expansions does not supply this lemma's
hypothesis. The finite and escaping pieces must together give an
invertible period matrix of the full de Rham rank. Also, a fractional
power becomes an integer after ramification; the proof uses its
strictly positive order, not an unsupported monodromy inference.

\subsection{A formal leading-coefficient criterion}\label{sec:stationaryballs}

For a polynomial phase, the remaining shrinking argument can be
replaced by a valuation calculation. The construction below includes
critical points that collide, remain multiple, or escape to infinity.
The key is to group them before making a Gaussian approximation.

Write $v$ for the valuation of a formal parameter $r$, normalized by
$v(r)=1$. The valuation of a polynomial is the least valuation of its
coefficients. We use the formal Newton--Puiseux theorem: finitely many
algebraic elements over $\C((r))$ belong to one field
$\C((r^{1/N}))$. Convergence is not required~\cite[Chapter~2]{Wall}.

\begin{lemma}[Polynomial separation of formal balls]\label{lem:ballseparation}
Let finitely many disjoint balls have centers $c_B$ in a Puiseux field
and radii specified by $v(x-c_B)\geq\rho_B$. After one common
ramification, put $\delta_B=r^{\rho_B}$. For any one ball $B$ and
integer $j\geq0$, there is a polynomial $q_{B,j}(x)$ over the Laurent
field whose substitution $x=c_C+\delta_C y$ has integral coefficients
for every ball $C$, and whose reduction is $y^j$ on $B$ and zero on
every other ball.
\end{lemma}
\begin{proof}
The finitely many valuations $v(c_B-c_C)$ give a finite tree of
clusters of centers. At a node with scale $r^d$, the coordinate
$u=(x-c)/r^d$ has distinct constant reductions $a_1,\ldots,a_s$
on its children. The Lagrange polynomial
\[
 e_i(u)=\prod_{k\ne i}\frac{u-a_k}{a_i-a_k}
\]
reduces to one on child $i$ and zero on every other child.

Start at the chosen leaf with $((x-c_B)/\delta_B)^j$. Induct upward.
Suppose a polynomial $h$ already has the required reductions on the
chosen child. On each of the other child balls its coefficients have
some finite lower valuation. The restriction of $e_i$ there has
strictly positive valuation. Consequently $e_i^M h$, for a sufficiently
large fixed integer $M$, has positive coefficient valuation on all
those children. On the chosen child $e_i^M=1+$ positive-order terms,
so the existing reductions are unchanged. There are finitely many
nodes, and this constructs the required polynomial. A node with just
one leaf requires no filter. This also explains why independently
normalized Lagrange polynomials, with no control of their growth on
other scales, would not suffice.
\end{proof}

\begin{theorem}[Formal stationary balls]\label{thm:stationaryballs}
Let $\varphi(r,x)\in\C((r))[x]$ have degree $D\geq2$, and let
$h(r)\in\C((r))\setminus\{0\}$. Put
\[
 M=\C((r))[x]/(\partial_x+\varphi_x)\C((r))[x],\qquad
 \nabla[q]=[q_r+(\varphi_r+h'/h)q].
\]
After a finite ramification, this module has a full frame of rank-one
formal exponential functionals. If $a_D$ is the leading coefficient
of $\varphi$, every zero-exponential functional in this frame takes
the seed $[1]$ to a series of valuation at least
\begin{equation}\label{eq:valuationbound}
 v(h)-\frac{v(a_D)}D.
\end{equation}
In particular, if $Dv(h)>v(a_D)$, the differential submodule generated
by $[1]$ has no horizontal quotient taking $[1]$ to the constant one.
\end{theorem}
\begin{proof}
The two twisted derivatives commute, so the parameter connection
is well defined on the quotient. A nonzero polynomial of degree $m$
has spatial twisted derivative of degree $m+D-1$. Reduction of
successive leading terms therefore gives a basis
$1,x,\ldots,x^{D-2}$ and dimension $D-1$.

Split $\varphi_x$, retaining root multiplicities. For a critical root
$\xi$, expand
\[
 \varphi(\xi+u)=c_\xi+\sum_{j=2}^D a_{\xi,j}u^j,
 \qquad
 \rho_\xi=\max_{2\leq j\leq D,\ a_{\xi,j}\ne0}
                    \frac{-v(a_{\xi,j})}{j}.
\]
Enlarge the ramification so that all these scales are integral, and
set $\delta_\xi=r^{\rho_\xi}$. The polynomial
\[
 \varphi(\xi+\delta_\xi y)-c_\xi
\]
has integral coefficients, no linear term, and at least one unit
coefficient. Its reduction $F_\xi(y)$ has degree $d_\xi\geq2$.
Here and below we may replace $r$ by a root of $r$ and renormalize
the valuation; every displayed strict inequality is unchanged.
As $a_{\xi,D}=a_D$,
\begin{equation}\label{eq:balllower}
 \rho_\xi\geq-v(a_D)/D.
\end{equation}

\emph{The balls partition all critical roots.}
If another critical root $\zeta$ satisfies
$v(\zeta-\xi)\geq\rho_\xi$, translate the integral scaled polynomial
by $(\zeta-\xi)/\delta_\xi$. All its coefficients remain integral,
and the coefficient of degree $d_\xi$ remains a unit: higher-degree
terms have positive valuation. The linear term vanishes because
$\zeta$ is critical. It follows from the definition of the maximum
that $\rho_\zeta=\rho_\xi$. Thus their balls are identical.
If two of the balls intersect, one contains the other; the same
argument makes them identical. Otherwise they are disjoint.

The scaled derivative has reduction $F_\xi'$ of degree $d_\xi-1$.
It therefore has exactly $d_\xi-1$ integral roots, with multiplicity.
Here is a direct check of that count: factor it in a Puiseux field.
An integral root contributes a linear factor to the reduction;
for a root of negative valuation, factor out that root and the
remaining factor reduces to a constant. Multiplicativity of the
coefficient valuation leaves a nonzero constant scalar. Hence the
degree of the reduction counts exactly the integral roots.
The ball $B$ thus contains $d_B-1$ critical roots. Since
$\deg\varphi_x=D-1$,
\begin{equation}\label{eq:ballrankcount}
 \sum_B(d_B-1)=D-1.
\end{equation}

\emph{A functional for every dimension in a ball.}
The fixed polynomial $F_B$ has twisted quotient
\[
 H_B=\C[y]/(\partial_y+F_B')\C[y]
\]
with basis $1,y,\ldots,y^{d_B-2}$. Let $\ell^0_{B,j}$ be its dual
basis. Write
\[
 \varphi(\xi_B+\delta_B y)=c_B+F_B(y)+E_B(r,y),
 \qquad E_B\in r\C[[r]][y].
\]
For a polynomial amplitude $q$ define
\begin{equation}\label{eq:ballfunctional}
 \ell_{B,j}(q)=h\delta_B e^{c_B}
   \ell^0_{B,j}\!\left(e^{E_B(r,y)}q(\xi_B+\delta_B y)\right).
\end{equation}
Each coefficient uses finitely many polynomial terms. The exponential
in this formula is purely formal. The product rule conjugates the
spatial twisted derivative to $\partial_y+F_B'$, so the functional
annihilates the defining relations of $M$.

We spell out the parameter check because it retains the moving volume.
At fixed $y$, put $w=\xi_B'+\delta_B'y$. Differentiating
\eqref{eq:ballfunctional}, then subtracting its value on
$q_r+(\varphi_r+h'/h)q$, leaves, inside the fixed functional,
\[
 e^{E_B}\left(wq_x+\frac{\delta_B'}{\delta_B}q+\varphi_xwq\right)
 =\frac1{\delta_B}(\partial_y+F_B')
                      \left(e^{E_B}wq\right).
\]
It is zero. Thus the functionals are horizontal, including the
Jacobian $\delta_B$ and normalization $h$.

\emph{Completeness, not just a rank count.}
Choose the polynomials $q_{B,j}$ of
Lemma~\ref{lem:ballseparation}, for $0\leq j<d_B-1$.
Divide each evaluation column by its scalar $h\delta_Be^{c_B}$.
The resulting matrix reduces to the identity: $E_B$ has positive
order and the reductions of the amplitudes are the chosen monomials
on one ball and zero elsewhere. The matrix is therefore invertible.
Leading-degree reduction gives $\dim M=D-1$, so
\eqref{eq:ballrankcount} accounts for the entire module.

Remove from $c_B$ its finite negative-power part $q_B(r)$.
The remaining exponential is an invertible formal power series,
up to a nonzero constant. Formula~\eqref{eq:ballfunctional} and the
invertible matrix give the asserted direct sum of rank-one factors
$e^{q_B}$ times powers of $r$. If $q_B=0$, its value on the seed
has valuation at least $v(h)+\rho_B$, proving
\eqref{eq:valuationbound}.

Finally, after ramification this is a semisimple differential module.
A quotient on the submodule generated by the seed extends to the
whole module. Nonzero exponential factors have no horizontal map
to the Laurent field, because the logarithmic derivative of a
Laurent series cannot be the derivative of a nonzero negative-power
exponent. On the zero part, every horizontal functional is a
constant combination of those in the full frame. If the bound is
strictly positive, their seed values cannot sum to one.
\end{proof}

The construction is a formal stationary-phase argument. Its classical
background includes Newton--Puiseux and the formal stationary-phase
theory of exponential connections; Sabbah's local Fourier--Laplace
formula is a primary reference for the latter~\cite{SabbahStationary}.
We do not identify that formula's scope with the particular proof
above or claim priority for this formulation. The proof here also
allows nonconvergent formal Laurent coefficients. It uses a valued
field, not an arbitrary coefficient ring.

\begin{corollary}[Leading-coefficient obstruction]\label{cor:leadingobstruction}
For a rational polynomial phase $\varphi(t,x)$ of degree $D\geq2$
and nonzero rational scalar normalization $h(t)$, let $v_\infty$
mean order at $\varepsilon=1/t=0$. If
\begin{equation}\label{eq:leadingcriterion}
 v_\infty(h^D/a_D)>0,
\end{equation}
its seed submodule cannot have a rational horizontal quotient with
seed value one. In particular any horizontal formal functional at
another point cannot take that seed to one.
\end{corollary}
\begin{proof}
Apply the theorem at infinity. A rational horizontal quotient remains
a quotient on extension to the formal Laurent field. A formal
functional with seed value one restricts to such a rational quotient:
its values on successive seed derivatives are $1,0,0,\ldots$.
\end{proof}

The inequality must be strict. For example
$\varphi=-x^2/(2t^2)$ and $h=1/t$ have
$v_\infty(h^2/a_2)=0$, and their normalized Gaussian period is a
nonzero constant. Thus the earlier adverse scaling example is also
an exact boundary test of the criterion.

\subsection{Closing the full rank-one $WT^2$ family}\label{sec:rankoneformal}

\begin{proposition}\label{prop:rankoneformal}
For $f\ne0$ and $\deg G\leq3$, the polynomial
\[
 P=W(a+bZ+cT+fT^2)+G(Z,T)
\]
is pure if and only if $b=0$ and $G=0$. Its pure members obey
$m>\deg Q$.
\end{proposition}
\begin{proof}
Put $q=1-bt$ and $A(T)=a+cT+fT^2$. Exact circular contraction
gives
\[
 J(p)=\frac1q\E_T\!\left[p(T)
             \exp\{tG(tA/q,T)\}\right],\qquad J(1)=\Phi_1.
\]
Gaussian integration by parts says that $J$ annihilates
$\partial_T+\varphi_T$; direct parameter differentiation makes it
horizontal for the polynomial phase and normalization
\[
 \varphi=tG(tA/q,T)-T^2/2,\qquad h=1/q.
\]
The phase has degree at least two, since it equals $-T^2/2$ at $t=0$.

Suppose first $b=0$ and write $G=\sum_{j\geq0}Z^jG_j(T)$.
If a term $A^jG_j$ attains the actual degree $D\geq2$ of the phase,
the leading coefficient is a polynomial in $t$ containing a
nonzero positive power. Distinct $j$ give distinct powers $t^{j+1}$,
so they cannot cancel identically. Its order at infinity is negative,
whereas $v_\infty(h)=0$. Criterion~\eqref{eq:leadingcriterion}
excludes purity. Because $\deg A=2$, every nonzero $G_j$ with
$j\geq1$ attains some degree at least two. The only remaining
possibility is $G=\alpha T+\beta$. Its exact unit series is
$\exp(\beta t+\alpha^2t^2/2)$, which is one only if $G=0$.

Suppose now $b\ne0$. Divide $G$ by $S=bZ+A(T)$ as a polynomial
in $T$:
\[
 G=R_0(Z)+TR_1(Z)+S(\alpha Z+\beta T+\gamma).
\]
The remainder has basis $1,T,Z,ZT,Z^2,Z^2T,Z^3$. At $Z=-A/b$
these monomials have distinct leading degrees $0,1,2,3,4,5,6$.
Thus if $V(T)=R(-A/b,T)$ has degree $n$, then
$\deg R(sA,T)\leq n$ for every scalar $s$; a constant $V$
forces $R$ to be that same constant. Write $g=tA/q$. The phase is
exactly
\[
 tR(g,T)+\alpha g^2+\beta gT+\gamma g-T^2/2.
\]
Since $g=-A/b+O(t^{-1})$, its first term is $tV+O(1)$, with
all these remainders of degree at most $n$.

If $\alpha\ne0$, the limiting remaining phase has degree four
with leading coefficient $\alpha f^2/b^2$. If $\alpha=0$,
$\beta\ne0$, it has degree three with leading coefficient
$-\beta f/b$. Otherwise its quadratic coefficient tends to
$-\gamma f/b-1/2$. A leading term from $tV$ has valuation $-1$;
a higher-degree nonzero limiting term has valuation zero.
The only flattened case has
$\alpha=\beta=0$, $\gamma=-b/(2f)$. Its exact quadratic
coefficient outside $tR$ is $-1/(2q)$, of valuation one.
Thus in every case the actual leading coefficient $a_D$ has
$v_\infty(a_D)\leq1$. This also covers $V$ constant or linear.
But $v_\infty(h)=1$ and $D\geq2$, so
$Dv_\infty(h)-v_\infty(a_D)\geq D-1>0$.
Purity is impossible by the same criterion.

Conversely $W A(T)$ has strictly negative circular weight.
For a monomial $Z^iW^jT^k$, a mixed moment can be nonzero only
when the number of $W$ factors balances the number of $Z$ factors;
in particular it vanishes for $m>\deg Q$.
\end{proof}

\subsection{A polynomial phase with logarithmic volume}\label{sec:puncturedpoly}

The puncture itself can be retained in the formal construction when
the phase is polynomial. This gives a stronger bound than the one for
$dx$: logarithmic volume does not change under a scaling of $x$.

\begin{theorem}\label{thm:puncturedpoly}
Let $K=\C((r))$, $\varphi\in K[x]$, and $h\in K\setminus\{0\}$.
On
\[
 H=K[x,x^{-1}]/(x\partial_x+x\varphi_x)K[x,x^{-1}],\qquad
 \nabla[q]=[q_r+(\varphi_r+h'/h)q],
\]
there is, after finite ramification, a full frame of rank-one formal
exponential functionals. Every zero-exponential functional in this
frame has seed value of valuation at least $v(h)$. Thus if $v(h)>0$,
the differential seed submodule has no horizontal quotient taking
$[1]$ to one. This includes a constant phase.
\end{theorem}
\begin{proof}
Write $D=\deg\varphi$, taking $D=0$ for a constant phase, including
zero. If $D\geq1$, the quotient has basis
$1,x,\ldots,x^{D-1}$. Indeed, on a negative monomial the lowest term
of $x\partial_x+x\varphi_x$ is its nonzero integer exponent times
that monomial. This eliminates all negative powers without introducing
lower ones. On polynomials, successive leading terms eliminate all
degrees at least $D$. Conversely an exact polynomial cannot come
from a Laurent polynomial with a negative lowest term, and a nonzero
polynomial preimage of degree $m$ has exact degree $m+D$.
This proves independence. For $D=0$, the quotient of $x\partial_x$
has the single basis element $1$.

Use the critical balls constructed in Theorem~\ref{thm:stationaryballs}.
At most one contains the origin. Every other ball has
$\rho_B>v(\xi_B)$, so $\gamma_B=\delta_B/\xi_B$ has positive
valuation. With the same $c_B,F_B,E_B$ and fixed functionals as before,
the local logarithmic functionals are
\begin{equation}\label{eq:puncturedaway}
 q\longmapsto h e^{c_B}\frac{\delta_B}{\xi_B}
  \ell^0_{B,j}\!\left(
      \frac{e^{E_B}q(\xi_B+\delta_B y)}{1+\gamma_B y}\right).
\end{equation}
The geometric denominator and all negative powers in $q$ expand
coefficientwise because $v(\gamma_B)>0$. Every coefficient is a
polynomial in $y$. The normalized logarithmic Jacobian is
$dx/x=\delta_B\,dy/(\xi_B+\delta_B y)$.

We retain the puncture in a separate central ball. If a critical
ball contains zero, write it as $v(x)\geq\rho_0$ and put
$\delta_0=r^{\rho_0}$. Translation within that ball gives
\[
 \varphi(\delta_0 y)=c_0+F_0(y)+E_0(r,y),\qquad
 c_0=\varphi(0),\quad E_0\in r\C[[r]][y],\quad F_0(0)=0.
\]
Here $d_0=\deg F_0\geq2$, and the ball contains $d_0-1$
critical roots with multiplicity. If no critical ball contains zero,
choose $\rho_0$ sufficiently large that the central ball is disjoint
from all critical balls and every nonconstant coefficient of
$\varphi(\delta_0 y)$ has positive valuation. Then $F_0=0$.
Such a choice exists: there are finitely many roots and coefficients.
This also handles $D\leq1$ without any critical balls.

For $d_0\geq2$, the fixed punctured quotient
\[
 H_0=\C[y,y^{-1}]/(y\partial_y+yF_0')\C[y,y^{-1}]
\]
has basis $1,y,\ldots,y^{d_0-1}$, by the same reduction argument.
For $F_0=0$, it has basis $1$ and dual functional the constant term.
For each element $\lambda_j$ of the dual basis define
\begin{equation}\label{eq:puncturedcentral}
 q\longmapsto h e^{c_0}\lambda_j
                   \!\left(e^{E_0}q(\delta_0 y)\right).
\end{equation}
There is no factor $\delta_0$: $dx/x=dy/y$. Each coefficient is
a Laurent polynomial, so the expression is well defined.

We check both identities, including parameter dependence. Spatial
exactness in the central ball is conjugation by $e^{E_0}$ of
$y\partial_y+yF_0'$. The derivative discrepancy is
$\delta_0'/\delta_0$ times the same exact expression applied to
$e^{E_0}q$. On an outside ball put $x=\xi_B+\delta_B y$ and
$w=\xi_B'+\delta_B'y$. After removing $h e^{c_B}$, the derivative
discrepancy in \eqref{eq:puncturedaway} is
\[
 e^{E_B}\left[
   \left(\frac{\delta_B'}x-\frac{\delta_B w}{x^2}\right)q
        +\frac{\delta_B w}x(q_x+\varphi_xq)\right]
 = (\partial_y+F_B')\left(e^{E_B}\frac wxq\right).
\]
It is annihilated by $\ell^0_{B,j}$. The identical product rule
proves spatial exactness there. Thus every displayed functional is
horizontal for the stated connection.

Completeness follows from polynomial separation of these disjoint
balls, not merely their number. For the central ball use the monomials
of the basis of $H_0$; for each outside ball use
$1,\ldots,y^{d_B-2}$. Lemma~\ref{lem:ballseparation} gives polynomial
amplitudes with exactly these reductions on their own ball and zero
on the others. Normalize the central columns by $h e^{c_0}$ and
the outside columns by $h e^{c_B}\delta_B/\xi_B$. Their matrix
reduces to the identity. If the central ball contains critical roots,
it contributes their multiplicity plus one; otherwise it contributes
one. Together with all outside roots this gives $D$ dimensions for
$D\geq1$, or one for $D=0$, the full quotient rank.

As in Theorem~\ref{thm:stationaryballs}, this frame is a direct sum
of rank-one exponential factors after ramification. For a zero factor,
\eqref{eq:puncturedcentral} has seed valuation at least $v(h)$;
\eqref{eq:puncturedaway} has valuation at least
$v(h)+\rho_B-v(\xi_B)>v(h)$. Semisimplicity now extends a quotient
of the differential seed submodule to the full module. Nonzero
exponential factors cannot map horizontally to the Laurent field.
A constant combination of the zero-factor seed values cannot be
one if $v(h)>0$.
\end{proof}

\begin{remark}[Why two irregular ends require another argument]
For $\varphi=r(x+x^{-1})$ and $h=1$, put $\theta=r\partial_r$.
Every logarithmic period satisfies
\[
 (\theta^2-4r^2)I=0.
\]
Indeed $y=x+x^{-1}$ and $z=x-x^{-1}$ obey
$(x\partial_x+x\varphi_x)z=y+rz^2$ and $y^2-z^2=4$.
These identities give the equation by integration by parts.
The quotient has rank two, and $[1],[y]$ are independent:
its basis is $x^{-1},1$ and $[y]=2[x^{-1}]$.

An explicit formal basis of scalar solutions is
\[
 I=\sum_{n\geq0}\frac{r^{2n}}{(n!)^2},\qquad
 J=I\log r-\sum_{n\geq1}\frac{H_n r^{2n}}{(n!)^2},
 \quad H_n=\sum_{j=1}^n\frac1j.
\]
The coefficient recurrence $n^2a_n=a_{n-1}$ proves the first
equation; $H_n-H_{n-1}=1/n$ proves the second.
Their $\theta$-Wronskian is one: its derivative is zero by the
equation and its constant term is one. Thus the logarithm is a full
rank-two phenomenon, with nontrivial unipotent monodromy that no
finite ramification removes. These are the familiar modified Bessel
solutions; see~\cite[10.25.1, 10.31.2]{DLMF}.
This example does not refute the desired logarithmic bound:
multiplication by $h=r$ makes both seed values tend to zero.
It does refute an extension of the semisimple polynomial-frame theorem
to all Laurent phases. The remaining proof must allow such extensions.
\end{remark}

The residue dimension at the puncture is essential. Dropping it would
give rank $D-1$ instead of $D$ and an incomplete frame. This theorem
still permits nonconvergent formal coefficients, but is a valued-field
statement. It does not cover a phase with both positive and negative
powers of $x$.

\begin{corollary}\label{cor:oneend}
Suppose $\varphi(t,x)$ is rational in $t$ and polynomial either in
$x$ or in $x^{-1}$. For the logarithmic-volume connection normalized
by $h=1/t$, no horizontal formal functional at another point can have
seed value one.
\end{corollary}
\begin{proof}
For a polynomial in $x$, apply Theorem~\ref{thm:puncturedpoly} at
$r=1/t=0$, where $v(h)=1$. The rational seed restriction is justified
as in Corollary~\ref{cor:leadingobstruction}. For a polynomial in
$x^{-1}$ use $y=1/x$; $dx/x=-dy/y$ only changes a constant factor.
\end{proof}

\subsection{A mixed chart with a complete formal proof}\label{sec:mixedformal}

One of the previously analytic exclusions admits a direct formal proof.
It is the whole family in the next proposition, not a generic choice
of its coefficients.

\begin{proposition}\label{prop:mixedformal}
For every $a\in\C$ and every $G\in\C[Z,T]$ of total degree at most
three, the polynomial
\[
 P=W(a+T+Z^2)+G(Z,T)
\]
has a nonzero positive pure Gaussian moment.
\end{proposition}
\begin{proof}
We construct the exact formal seed functional and apply
Corollary~\ref{cor:leadingobstruction}. No analytic integral is needed.

\emph{The formal source.}
The circular fixed point is $g=t(a+T+g^2)$, with determinant
$D=1-2tg$. Make the formal Gaussian change
\[
 T=Y-t^2(a+Y)^2,\qquad z=t(a+Y).
\]
It has $g=z$ and $dT/dY=D$. For any change $T(Y)=Y+O(t)$ with
polynomial formal coefficients, the identity
\[
 \E_T[F(T)]=\E_Y\!\left[F(T(Y))T'(Y)
             e^{(Y^2-T(Y)^2)/2}\right]
\]
holds coefficientwise. To prove it, use the homotopy
$T_u=Y+u(T-Y)$: its derivative is the Gaussian expectation of
$(\partial_Y-Y)\{F(T_u)\dot T_u e^{(Y^2-T_u^2)/2}\}$, which is zero.
Thus the determinant cancels and gives
\[
 \Phi_1=\E_Y[e^\Psi],\qquad
 \Psi=tG(z,Y-z^2)+Yz^2-z^4/2.
\]

Define
\[
 F=z/t-a-z^2,\qquad \varphi(t,z)=tG(z,F)-F^2/2,
\]
and let
\[
 M=K[z]/(\partial_z+\varphi_z)K[z],\qquad
 \nabla[q]=[q_t+(\varphi_t-1/t)q],\qquad K=\C(t).
\]
These are vector-space quotients. The formal map
$J(q)=\E_Y[q(t(a+Y))e^\Psi]$ annihilates
$(\partial_z+\varphi_z)K[z]$ by Gaussian integration by parts.
Differentiating at fixed $Y$ differs from applying $\nabla$ by
$J((\partial_z+\varphi_z)(zq/t))=0$.
So $J$ is horizontal and $J(1)=\Phi_1$. The term $-1/t$ is the
normalization $dz/t=dY$; removing it changes the problem.
Leading-degree reduction gives $\dim_K M=\deg_z\varphi-1$.

\emph{The exhaustive degree split.}
Put
\[
 V(z)=G(z,-a-z^2),\quad L(z)=G_T(z,-a-z^2),\quad
 B_1=z^2G_{TT}(z,-a-z^2)/2+z(a+z^2).
\]
If $s=[T^3]G$, direct Taylor expansion gives
\begin{equation}\label{eq:mixedphi}
 \varphi=tV+B_0+B_1/t+B_2/t^2,
 \quad B_0=zL-(a+z^2)^2/2,\quad B_2=sz^3-z^2/2.
\end{equation}
To check every degree drop, write $v_1=[ZT^2]G$,
$v_0=[T^2]G$ and $u_2=[Z^2T]G$. Then
\[
 [z^6]V=-s,\qquad [z^5]V=v_1,\qquad
 [z^4]V=v_0-u_2-3as.
\]
If $n=\deg V\geq4$, the degree of $\varphi$ is $n$ and each $B_j$
has degree at most $n$. If $n\leq3$, or $V$ is constant, these three
identities force $s=v_1=0$ and $u_2=v_0$. Hence
$B_0=-z^4/2+O(z^3)$, $\deg B_1\leq3$ and $B_2=-z^2/2$.
In that case $\deg_z\varphi=4$. There is no omitted higher-degree
term tending to zero at infinity.

If $n\geq4$, the leading coefficient of $\varphi$ has a nonzero
$tV_n$ term and therefore has order $-1$ at infinity. If $n\leq3$
or $V$ is constant, its leading coefficient is exactly $-1/2$ and
has order zero. In both cases $h=1/t$ has order one, so
\[
 v_\infty(h^D/a_D)>0.
\]
The leading-coefficient obstruction excludes $J(1)=1$. The degree
calculation was exhaustive, so this proves the proposition even at
multiple critical points and degree drops. The full formal frame,
including every escaping dimension, is supplied by
Theorem~\ref{thm:stationaryballs} rather than a new case-by-case
saddle calculation.
\end{proof}

\subsection{An exhaustive split of the affine-linear coordinate}\label{sec:wsplit}

Write $h=a+bZ+cT+dZ^2+eZT+fT^2$ and
$\Delta=e^2-4df$. The following split is exhaustive and its branches
are disjoint. The last column distinguishes a proof given here from
the proposition proving that branch. Every row is now proved here.

\begin{center}\small
\begin{tabular}{@{}p{0.27\textwidth}p{0.38\textwidth}p{0.25\textwidth}@{}}
\toprule
Coefficient condition & Reduction & Status in this draft\\
\midrule
$\Delta\ne0$ & Laurent phase on $\C^*$ with logarithmic volume &
Proposition~\ref{prop:nondeglog}\\[3pt]
$\Delta=0$, $f\ne0$ & $h=a+bZ+cT+fT^2$ after a null rotation &
Proposition~\ref{prop:rankoneformal}\\[3pt]
$f=e=0$, $d\ne0$, $c\ne0$ &
$W(a+T+Z^2)+G$ after normalization &
Proposition~\ref{prop:mixedformal}\\[3pt]
$f=e=0$, $d\ne0$, $c=0$ & $Wh(Z)+G$ &
Proposition~\ref{prop:height}\\[3pt]
$d=e=f=0$, $c\ne0$ & $W(a+bZ+cT)+G$ &
Corollary~\ref{cor:linearwt}\\[3pt]
$d=e=f=c=0$ & $Wh(Z)+G$ &
Proposition~\ref{prop:height}\\
\bottomrule
\end{tabular}
\end{center}

Here is the algebra behind the reductions. The null rotation
\[
 Z=Z',\quad T=T'+\lambda Z',\quad
 W=W'-\lambda T'-\lambda^2Z'/2
\]
preserves $2ZW+T^2$. If $f\ne0$, choose $\lambda=-e/(2f)$.
When $\Delta=0$ the quadratic part of $h$ becomes just $fT'^2$.
If $f=e=0$ and $c\ne0$, choose $\lambda=-b/c$ to remove the
linear $Z$ term without changing $dZ^2$. If also $d\ne0$, the
Gaussian scaling $Z=\kappa Z'$, $W=W'/\kappa$, followed by a
nonzero scaling of $P$, makes the $T$ and $Z^2$ coefficients both
one; take $\kappa^2=c/d$. All additional terms enter a polynomial
$G$ of degree at most three. If $\Delta=0$ and $f=0$, then $e=0$,
which proves that the table omits no case.

\subsection{A curve-length estimate}

\begin{lemma}\label{lem:curvelength}
A real algebraic curve of degree at most $D$ has length at most
$\pi D\rho$ in a disk of radius $\rho$. Consequently its logarithmic
length in $a\leq|x|\leq b$ is at most a constant times
$D(1+\log(b/a))$.
\end{lemma}
\begin{proof}
Average the number of intersections with lines over directions and
apply the one-dimensional coarea formula. Outside exceptional line
components the number is at most $D$; the components themselves
satisfy the same estimate. In each dyadic annulus
$\rho\leq|x|\leq2\rho$, divide the disk estimate by $\rho$ to
bound $\int|dx|/|x|$ by $2\pi D$. Sum the dyadic annuli.
\end{proof}

\subsection{The logarithmic branch, including its removed component}

For $\Delta\ne0$, $f\ne0$, choose $\eta^2=\Delta$ and set
\[
 \rho=\frac{e+\eta}{2f},\quad \sigma=\frac{e-\eta}{2f},\quad
 u=\frac{\rho c-b+1/t}{\eta},\quad
 v=\frac{b-1/t-\sigma c}{\eta},\quad k=uv-a/f.
\]
On the circular constraint $z=th(z,T)$, put
$x=T+\rho z+u$, so that $T+\sigma z+v=k/x$. Then
\[
 z=\frac f\eta(x-k/x-u+v),\quad T=x-u-\rho z,
 \qquad
 \frac{dT}{1-th_Z}=\frac{dx}{t\eta x}.
\]
The coefficient of $t^{-2}$ in $k$ is $-1/\eta^2$, so the conic
is generically smooth. The exponent $tG(z,T)-T^2/2$ is Laurent
polynomial in $x$. On the formal Gaussian branch,
$x=1/(\eta t)+O(1)$, which makes every Laurent power legitimate.

If $f=0$, then $e\ne0$. Put $z_0=-c/e$ and $x=z-z_0$. With
$A(z)=a+bz+dz^2$, the constraint gives
\[
 T=\frac{z/t-A(z)}{ex},\qquad
 \frac{dT}{1-th_Z}=\frac{dx}{tex}.
\]
This localization is valid if $c\ne0$, or if $c=0$ and $a\ne0$.
If $a=c=0$, the chosen fixed point is instead $g=0$ identically;
inverting $g$ would discard the very branch being tested. In that
case the exact unit series is
\[
 \E_T\left[\frac{e^{tG(0,T)}}{1-t(b+eT)}\right].
\]
Use $x=T-(1/t-b)/e$, now a valid Laurent unit. Its volume is
$-dx/(tex)$ and its exponent is polynomial in $x$. These formulas
cover both the smooth conic and the removed-component branch.
We next connect these coordinates to the formal Gaussian functional.
This will close every one-ended phase without an analytic estimate.

\begin{lemma}[Formal logarithmic transport]\label{lem:logtransport}
In each of the three charts above, write $X=X(t,T)$ for the formal
coordinate, $D$ for the displayed circular denominator, and $B$ for
$tG(g,T)$ or $tG(0,T)$ on the removed component. Put
$\kappa=\eta,e,-e$, respectively. Then
\[
 J(q)=\E_T\!\left[\frac{q(X)e^B}{D}\right]
\]
is horizontal on the logarithmic twisted quotient with phase
$\varphi$ and normalization $h=1/t$, and $J(1)=\Phi_1$.
\end{lemma}
\begin{proof}
All Laurent powers of $X$ are legitimate by the unit checks above.
The coordinate identities give
\[
 X_T=t\kappa X/D,\qquad
 \varphi(t,X)=B-T^2/2.
\]
Consequently Gaussian integration by parts gives
$J((x\partial_x+x\varphi_x)q)=0$: the corresponding integrand,
multiplied by $t\kappa$, is
$(\partial_T-T)\{q(X)e^B\}$.

For parameter horizontality set $w=1/D$ and $R=X_t/X_T$.
Differentiating $w=X_T/(t\kappa X)$ gives the exact identity
\[
 w_t+w/t=\partial_T(wR).
\]
Subtracting $J(q_t+(\varphi_t-1/t)q)$ from $\partial_tJ(q)$ now
gives the Gaussian expectation of
$(\partial_T-T)\{wR q(X)e^B\}$, hence zero. These are formal
coefficientwise identities, without an integral-convergence hypothesis.
The seed equality is the circular contraction formula.
\end{proof}

\begin{proposition}\label{prop:oneendgaussian}
Every chart above whose actual Laurent phase has at most one
irregular end is not pure. The removed-component branch
$a=c=f=0$, $e\ne0$, is therefore excluded for every polynomial $G$,
without a degree bound.
\end{proposition}
\begin{proof}
At most one irregular end means that $\varphi$ is polynomial in
$x$ or in $x^{-1}$, including the constant case. Purity would give
the horizontal functional in Lemma~\ref{lem:logtransport} seed value
one, contrary to Corollary~\ref{cor:oneend}. The removed-component
phase is always polynomial in $x$, irrespective of $\deg G$.
\end{proof}

\subsection{Resolving a Laurent phase}\label{sec:laurentresolution}

The logarithmic estimate needs two ingredients: a complete description
of the regular part, and a bound for the distinguished differential
$d\log x$. The first is a nearby-cycle calculation. The second is a
calculation on annuli. Keeping these ingredients separate prevents a
perfect period pairing from being mistaken for a growth estimate.

Throughout this subsection the coefficients are rational functions of
the parameter. In particular their Puiseux expansions converge. The
nonconvergent formal statements proved earlier have a different scope.

The two-ended twisted quotient is finite dimensional without a
genericity assumption. If the actual Laurent orders are $p,q\geq1$,
the classes $x^{-q},\ldots,x^{p-1}$ are a basis. The top term of
$x\varphi_x$ eliminates every exponent at least $p$, without crossing
below $-q$; the bottom term eliminates every exponent below $-q$,
without crossing above $p-1$. An exact expression supported in this
interval would require its preimage to have highest exponent at most
$-1$ and lowest exponent at least zero. This proves independence
and rank $p+q$. Coefficient specializations use their actual orders.

\begin{lemma}[A marked model]\label{lem:markedmodel}
Let $\varphi(r,x)\in\C(r)[x,x^{-1}]$ be nonconstant in $x$.
After a finite ramification of $r$ there is a smooth surface proper
over a parameter disc, unchanged over its punctured disc, with these
properties:
\begin{enumerate}
\item Its generic fibre is $\mathbb P^1$. Its central fibre is a
reduced normal-crossing curve, with local equations $r=u$ or $r=uv$.
\item The phase extends to a morphism to $\mathbb P^1$. Its zero
and pole divisor and the divisor of $x$ have normal crossings with
the central fibre.
\item The horizontal sections $x=0,\infty$ are disjoint and meet
smooth points of the central fibre. Locally $x=u^a v^b U$, where
$a,b$ are integers and $U$ is holomorphic and nonvanishing.
\end{enumerate}
Additional markings used in the construction are not additional
punctures of the exponential connection.
\end{lemma}
\begin{proof}
Mark $0,\infty$ and every point of the generic fibres
$\varphi^{-1}(0),\varphi^{-1}(1),\varphi^{-1}(\infty)$.
Retain multiplicities in the divisors, but identify identical sections.
Convergent Newton--Puiseux makes the finitely many marked branches
meromorphic sections after one finite ramification. Start with
$\mathbb P^1$ over that disc and blow up collisions of distinct
sections. This terminates: the contact order of each pair is a finite
integer and decreases on blowing up their common point.

Every centre is a smooth point of a single central component. Indeed
a section cannot pass through a node $uv=r$: two positive integral
orders cannot have sum one. Thus each new exceptional component has
multiplicity one and all nodes still have the equation $uv=r$.
The final marked sections are disjoint and miss the nodes.

The divisor of the meromorphic function $\varphi$ is supported on
the central components and its marked zero and pole sections. If a
point were an indeterminacy, its local divisor would have a zero
branch and a pole branch meeting there. After cancelling common
factors its expression is a unit times $u^a/v^b$, $a,b>0$.
The equation $\varphi=1$ then has a horizontal branch through that
point. At a node this contradicts the section property; at a smooth
point it contradicts the separation from the zero or pole section.
These are all possible intersections. Hence $\varphi$ is a morphism.
The asserted normal crossings and the expression for $x$ follow
from its divisor: horizontally $x$ has just the two marked sections.
Only those original sections are removed from the generic fibre;
marking a zero or a point with phase value one does not remove it.
\end{proof}

The marking of phase value one is necessary. For $\varphi=r/x^2$,
the zero and pole sections are already disjoint, but the map is
undefined at $r=x=0$. Its value-one branches $x=\pm\sqrt r$
force the cover $r=s^2$ and the chart $x=sy$, where the phase is
$1/y^2$. The marked construction therefore detects a transition
that merely separating zeros and poles would miss.


We use the following precise comparison facts. For a meromorphic
connection on a punctured disc, its moderate nearby cycles describe
its formal regular summand, with monodromy. Nearby cycles commute
with proper direct image. For a rank-one exponential on a surface
with normal crossings, a component on which the phase has a vertical
pole contributes no moderate nearby cycles. At a node between a
finite-phase component and a pole component, the surviving object is
the exponential connection on the finite component, localized at the
node; it is not an extra point class. At a node between two finite
components, the local complex is the usual nearby annulus complex,
including its monodromy. These are standard comparison and
normal-crossing facts, not new theorems here; see
Roucairol~\cite[Section~1.2 and Lemma~1.4]{RoucairolFormal} and
Sabbah~\cite[Section~1.b, especially Proposition~1.5]{SabbahStationary}.

For clarity, their local exponential models are a holomorphic unit
(finite phase), $\exp(u^{-p})$ at a finite--pole boundary, and
$\exp(u^{-p}v^{-q})$, $p,q>0$, when both components are poles.
An analytic unit in a pole monomial is absorbed in a coordinate.
This may multiply the equation of the base by a nonvanishing unit;
locally that change does not change the vanishing assertion or the
nearby-cycle object up to the usual unit identification. In the
$f^s$ description the identification is multiplication by
$\exp(s\log U)$, with a local analytic choice of $\log U$.
One keeps the actual gluing maps and monodromy, rather than just
adding the ranks of the components.

\begin{lemma}[Lifting the regular part]\label{lem:regularlifting}
For the marked model, the complete nearby-cycle space of the
exponential direct image is represented, on a sufficiently small
parameter sector, by flat rapid-decay chains made from the following
pieces:
\begin{enumerate}
\item bounded deformations of finitely many paths and cycles in
the finite-phase central components, away from their nodes and poles;
\item a bounded number of bridges and circles in the necks
$uv=r$ between finite-phase components;
\item descending phase tails at irregular ends and circles about
regular punctures.
\end{enumerate}
The coefficients of each chosen flat chain are fixed constants.
Their nearby-cycle classes exhaust the formal regular part, with
its actual gluing maps and monodromy.
\end{lemma}
\begin{proof}
Apply proper nearby-cycle comparison to the compactified surface,
with the exponential connection localized at $r=0$ and at the two
original punctures. The normal-crossing rules just stated kill the
interiors of the pole components. Along a finite component they leave
its ordinary exponential connection, with localization at its pole
ends and at any original regular puncture. A finite--pole node gives
that localized end. A finite--finite node retains the nearby annulus.
Thus no pole component is silently counted as an additional regular
class, and no annulus or puncture residue is discarded.

Here is a chain description of this calculation. Remove small discs
about the finitely many nodes and punctures in each finite component.
Triangulate the remaining compact surfaces, keeping the boundary
circles. At an irregular boundary attach the finitely many descending
sectors. At a regular puncture keep the circle. The one-dimensional
exponential period comparison identifies the resulting relative
chain complexes with the component de Rham complexes; for a constant
phase this is ordinary homology. At each finite--finite node glue by
the actual annulus chain complex, with both its degree-zero and
degree-one terms. Taking the total complex is Mayer--Vietoris for
the full nearby-cycle complex above. It computes its maps and
monodromy, not merely its Euler characteristic. In particular its
classes exhaust the formal regular summand by proper comparison.

Choose a finite basis of this finite-dimensional homology. Smooth
pieces and their boundary points lift by the implicit function
theorem; an annulus is lifted by $uv=r$. Descending ends are continued
by the inverse phase coordinate. On a fixed small sector these choices
are isotopies, so the homology coefficients are constants. No Morse
condition on the phase is used. Degenerate critical points stay
inside the triangulated pieces, and a constant phase on a component
is allowed.

For the tail continuation, choose the ray of $r$ so that no unbounded
critical value lies on the negative real phase axis for small $r$.
There are finitely many excluded directions: critical values are
convergent Puiseux functions. Bounded critical values fit in a fixed
disc. From a sufficiently negative phase value every chosen inverse
branch therefore continues to $-\infty$ without meeting a critical
value. In the two-ended Laurent case it ends at $x=0$ or $\infty$.
The same choice can avoid the finitely many limiting Stokes directions
of the parameter connection.

This constructs flat lifts whose nearby-cycle classes exhaust the
space computed by proper comparison. No quantitative estimate has
been used in that completeness calculation.
\end{proof}

\begin{lemma}[The logarithmic bound on the lifted pieces]
\label{lem:logpieces}
For the chains of Lemma~\ref{lem:regularlifting}, the period of
$e^{\varphi}d\log x$ is $O(1+|\log r|)$ along the chosen ray.
The bound is uniform for a fixed finite collection of generators.
With a fixed Laurent amplitude the same construction has at most
an additional algebraic power of $|r|^{-1}$.
\end{lemma}
\begin{proof}
On a compact finite-phase piece the phase and the coefficients of
$d\log x$ are bounded, so its integral is bounded. In a finite--finite
neck write
\[
 r=uv,\qquad x=u^a v^b U(u,v).
\]
At fixed $r$,
\[
 d\log x=\left(a-b+\frac{uU_u-vU_v}{U}\right)\frac{du}{u}.
\]
The coefficient is uniformly bounded in a small bidisc. A radial
bridge across $|r|/\eta\leq|u|\leq\eta$, with bounded angular
detours, has logarithmic length $O(1+|\log r|)$. Circles have bounded
logarithmic length. The phase is holomorphic on this bidisc and is
uniformly bounded there. This treats every finite--finite neck.

Near an irregular end of a finite component, use a local root of
the phase unit to write $\varphi=w^{-p}$, $p\geq1$. At a
finite--pole node one can simultaneously write $r=wv$ after changing
the other coordinate by a unit. On $|v|,|w|<\eta$, the same
calculation bounds $d\log x$ by a constant times $|dw|/|w|$.
Along a descending phase ray, $s=-\varphi>0$, its integral is bounded
by a constant times
\[
 \int_{s_0}^{\infty}e^{-s}\frac{ds}{s}.
\]
If the node chart ends, truncate it first at $|w|=|r|^\alpha$,
where $0<\alpha<1$. Its phase is already of size
$-|r|^{-\alpha p}$. Continue that inverse branch along the negative
phase axis as in Lemma~\ref{lem:regularlifting}.

We give a uniform bound for that continuation, rather than declaring
it negligible from a local saddle formula. Choose $M$ large enough
to include all coefficient transition scales in
$|r|^M\leq|x|\leq|r|^{-M}$. Such $M$ exists because the coefficients
are rational and there are finitely many Laurent powers. In this
annulus each descending path lies on $\operatorname{Im}\varphi=0$,
a real algebraic curve of bounded degree after clearing its Laurent
denominator. The curve-length bound used in
Lemma~\ref{lem:curvelength} gives
$\int|dx|/|x|=O(1+|\log r|)$, by dyadic annuli. Along the entire
tail the exponential is at most
$\exp(-|r|^{-\alpha p})$. Outside that annulus the leading Laurent
monomial dominates uniformly, so
$|d\log x/d\varphi|\leq C/|\varphi|$ and its descending integral
is bounded by the same exponential tail integral. This proves the
required uniform tail bound. A fixed Laurent amplitude contributes
only a power of $|r|^{-1}$ on the annulus and algebraic growth in
$s$ outside; the exponential still dominates it.

At a regular puncture use a small fixed circle in its smooth marked
chart. Both its phase and logarithmic length are bounded. There are
finitely many pieces and their chain coefficients are fixed.
Adding their bounds proves the lemma.
\end{proof}

\begin{theorem}[Logarithmic seed bound]\label{thm:laurentlog}
Let $\varphi(t,x)\in\C(t)[x,x^{-1}]$. On a generic ray at
$t=\infty$, the zero-exponential part of every scalar flat period
of $e^{\varphi(t,x)}dx/x$ is $O(1+\log|t|)$. This includes all
critical degeneracies, coincident actions and regular punctures.
Consequently, for a rational normalization $h(t)$ with
$v_\infty(h)>0$, its differential seed submodule admits no rational
horizontal quotient taking the seed to one. The same exclusion
holds for a horizontal formal Gaussian functional at another point.
\end{theorem}
\begin{proof}
If the phase is independent of $x$, the sole residue period is a
constant times $e^{\varphi(t)}$. Its zero-exponential part is bounded
when $\varphi$ is regular at infinity, and absent otherwise.
For a nonconstant phase put $r=1/t$ and use
Lemmas~\ref{lem:markedmodel}--\ref{lem:logpieces}. The construction
includes regular puncture circles when there is just one irregular end.
Finite ramification changes $\log|r|$ only by a constant factor.
The lifted chains have at most algebraic growth on every Laurent
amplitude by Lemma~\ref{lem:logpieces}. On the chosen generic ray,
a full period vector with such growth has no exponentially growing
factor: after formal exponential decomposition a largest growing
factor could not cancel on every basis amplitude. The other nonzero
factors are exponentially decaying there. The nearby-cycle basis
therefore gives a full basis of regular functionals modulo those
decaying terms, with constant coefficients in the flat basis. Every
one of their seed periods has the asserted logarithmic bound. Thus
their zero-exponential seed
expansions have no negative algebraic powers. Multiplication by $h$
makes all of them tend to zero.

For a constant phase the rank-one period already gives the same
contradiction directly. For a nonconstant phase, a quotient with
normalized seed one would extend locally to a
horizontal functional on the full module, by the seed restriction
and analytic extension argument of Lemma~\ref{lem:seedfunctional}.
Its comparison hypotheses hold on
$(\mathbb A^1\setminus S)\times\mathbb G_m$: remove the finitely
many parameter poles, zeros of $h$, and comparison exceptions.
The base projection and nonconstant phase are regular and
algebraically independent there.
Its scalar seed is constant. Distinct formal exponential factors
separate, so its zero-exponential seed is also one. This contradicts
the bound after normalization. A formal Gaussian functional with
unit seed first restricts to the same rational seed quotient, as in
the first part of Lemma~\ref{lem:seedfunctional}.
\end{proof}

The comparison theorems supply completeness here. The bounds come
from the logarithmic differential on the marked model and the explicit
descending-tail estimate. Neither part follows from a finite symbolic
calculation. This proof uses convergence of the rational parameter
family; it makes no assertion for arbitrary nonconvergent Laurent
families or arbitrary coefficient rings.


\begin{proposition}\label{prop:nondeglog}
If $e^2-4df\ne0$, no polynomial $G$, of any degree, makes
$P=Wh(Z,T)+G(Z,T)$ pure.
\end{proposition}
\begin{proof}
The smooth-conic, linear-conic and removed-component charts are
exhaustive. Their Laurent-unit checks and
Lemma~\ref{lem:logtransport} give the formal Gaussian functional
with seed value $\Phi_1$ and normalization $1/t$.
Theorem~\ref{thm:laurentlog} excludes $\Phi_1=1$ in every chart,
using its actual Laurent orders. No image-kernel calculation is used.
\end{proof}

\begin{proof}[Proof of Proposition~\ref{prop:wlinear}]
Use the disjoint case table in Section~\ref{sec:wsplit}.
Proposition~\ref{prop:nondeglog} excludes its first row.
Proposition~\ref{prop:rankoneformal} treats the second and leaves
only a one-sign product. Proposition~\ref{prop:mixedformal} excludes
the third. Proposition~\ref{prop:height} proves the fourth and sixth;
its pure cases are one-sign phases or the shear family.
Corollary~\ref{cor:linearwt} treats the fifth and leaves a quadratic.
These rows exhaust all coefficients, including their zero faces.

The one-sign and ternary quadratic cases have cutoff $m>\deg Q$,
and the shear has cutoff $m>3\deg Q$. The null rotations,
Gaussian coordinate scalings and nonzero scalar multiples of $P$
used in the reductions preserve the threshold and multiplier degree.
Thus every pure member of the full affine-linear family satisfies
$m>3\deg Q$.
\end{proof}

\subsection{Adversarial tests for the analytic argument}

The following exact examples prevent several tempting shortcuts.
For positive real $t$,
\[
 \frac1t\int_{\mathbb R}e^{-x^2/(2t^k)}\,dx
   =\sqrt{2\pi}\,t^{k/2-1}.
\]
The unique critical point is bounded and has zero action for every
$k$, yet the normalized period decays for $k=1$, is constant for
$k=2$, and grows for $k=3$. A critical-point count or a zero-action
label alone gives no decay estimate.

For the fixed phase $tx^k$, $k>2$, the critical point is degenerate
and the decay scale is $t^{-1/k}$, not a Gaussian Hessian scale.
The fixed quotient in Theorem~\ref{thm:stationaryballs}
retains its full multiplicity $k-1$.
For $tV_nx^n-x^4/2$, $1\leq n\leq3$, the $n-1$ bounded and
$4-n$ escaping critical dimensions total three. Counting only
bounded points loses actual rank; the stationary-ball construction
retains all roots and tests its full matrix using polynomial separation.
Finally, $t^{-3/2}=r^3$ under $r=t^{-1/2}$: a claim about a
noninteger exponent before ramification must not be reused after
ramification. Lemma~\ref{lem:positiveformal} uses positive seed order
and an actual full frame instead.

\subsection{One exclusion criterion for several geometric cases}
\label{sec:gradientexclusion}

The smooth case does not need a list of critical-point formulas.
The same argument also treats many singular leading forms. Its
hypotheses must hold in every complex direction. A bound only on
the original real integration space would not suffice.

We state the criterion for a polynomial exponential amplitude, then
specialize it to the Gaussian. Norms are Hermitian norms on complex
affine space; the Gaussian quadratic remains holomorphic.

\begin{theorem}[A gradient criterion]\label{thm:gradientcriterion}
Let $P,q\in\C[x_1,\ldots,x_n]$, with $P$ nonconstant and
$\deg q\leq J$, where $J\geq1$. Suppose that outside a compact set
\begin{equation}\label{eq:gradientconditions}
 \|\nabla P(x)\|\geq c\|x\|^{J-1},\qquad
 \frac{\|\operatorname{Hess}P(x)\|}{\|\nabla P(x)\|^2}\leq C
\end{equation}
for some $c>0$ and $C<\infty$. For every fixed polynomial $A$,
no nonzero flat period of
\[
 A(x)e^{tP(x)-q(x)}\,dx_1\cdots dx_n
\]
is a finite exponential-polynomial sum
$\sum_\lambda e^{\lambda t}p_\lambda(t)$ with polynomial
$p_\lambda$. In particular it cannot be a nonzero constant.
\end{theorem}
\begin{proof}
We first justify the full cycle space and the tails. The gradient is
nonzero outside a compact set, so the affine critical set is compact
and algebraic, hence finite. The gradient bound is stronger than the
condition excluding vanishing cycles at infinity. The finite-critical
thimble decomposition therefore applies, with the whole local Milnor
homology at each isolated critical point, not just its Morse
approximation~\cite[Part~I, Sections~1--4]{Pham}.
When several critical points have the same value, take disjoint
Milnor balls and apply excision to their complement. This gives the
direct sum of their full local homology before any values are grouped.

Here the Gaussian or polynomial amplitude can be kept on these fixed
$P$-cycles. Choose a generic argument $\theta$ of $t$. Outside the
local critical neighbourhoods, lift a descending phase ray by
\[
 P(x)=\lambda-e^{-i\theta}s,\qquad
 \frac{dx}{ds}=-e^{-i\theta}
                    \frac{\overline{\nabla P(x)}}{\|\nabla P(x)\|^2}.
\]
For large $s$ the lift is outside the compact set, because its phase
value is unbounded. The first condition in
\eqref{eq:gradientconditions} gives
$\|x\|^J\leq C_1(1+s)$. The real differential of this vector field
has norm at most a constant times the second ratio in
\eqref{eq:gradientconditions}, hence is bounded. Transported
$(n-1)$-cycle volumes grow at most as $\exp(C_2s)$; the radial speed
is bounded as well. Also
$|A(x)e^{-q(x)}|\leq C_3(1+s)^M\exp(C_4s)$ for some $M$.
Thus $\exp(-|t|s)$ dominates the complete tail, uniformly for large
$|t|$ in a smaller sector. The same estimate applies to every fixed
polynomial amplitude.

We must also check that these are a full basis for the actual phase,
rather than selected admissible cycles. Put
\[
 F_\sigma=P-\sigma q/t,\qquad 0\leq\sigma\leq1.
\]
Since $\|\nabla q\|\leq C_5\|\nabla P\|$ outside the same compact
set, all large $|t|$ give
\[
 \|\nabla F_\sigma\|\geq\tfrac12\|\nabla P\|.
\]
The ratio involving the Hessian stays uniformly bounded too. Indeed
$\|\operatorname{Hess}q\|=O(\|x\|^{J-2})$ for $J\geq2$, while
that Hessian is zero for $J=1$.
All critical points and their values therefore remain in fixed compact
sets. On the preimage of a phase half-plane beyond those values, the
parameter field
\[
 W_\sigma=\frac{q}{t}
       \frac{\overline{\nabla F_\sigma}}{\|\nabla F_\sigma\|^2}
\]
satisfies $\partial_\sigma F_\sigma+dF_\sigma(W_\sigma)=0$ and
$\|W_\sigma\|\leq C_6(1+\|x\|)/|t|$. Its flow is complete for
$0\leq\sigma\leq1$. Extend it by a cutoff away from that half-plane.
It gives an isotopy of the ambient relative pairs defining the full
rapid-decay homology. This explicitly prevents an infinity class
from appearing during the deformation.

The fixed $P$-thimbles have uniformly descending tails for every
$F_\sigma$, by the preceding estimates. Truncate them beyond a common
large phase value. Their boundaries stay in the descending half-plane
throughout the isotopy, so they define its flat relative classes.
They are a full basis at $\sigma=0$ and remain one at $\sigma=1$.
Finite critical points may merge, and their critical values may
coincide; the argument uses the full relative pairs and does not
require a Morse family. The same cycles are flat as $t$ varies in
the chosen sector.

These improper chains and their homotopies are admissible for the
period pairing. One can integrate first on truncated chains and then
pass to the limit: the tail bound makes every boundary term tend to
zero. Indeed the differential of $W_\sigma$ grows at most as
$C(1+\|x\|^J)/|t|$; along the controlled tails this is
$C'(1+s)/|t|$. The homotopy volumes have at most this exponential
growth, still dominated by the descending factor for all large $|t|$.
Thus Stokes and differentiation under the integral hold throughout;
the relative-pair calculation does not depend on a convergence claim
for the original real Gaussian integral.

It remains to determine the local asymptotics at a finite critical
point with value $\lambda$. For an analytic amplitude germ, the
vanishing-cycle density is a finite sum of terms
$s^{\alpha+k}(\log s)^j$, with analytic coefficient series and
$\alpha>-1$. Its Laplace asymptotics consequently have the form
\[
 e^{\lambda t}\sum_{\alpha,k,j}
       c_{\alpha,k,j}\,t^{-\alpha-k-1}(\log t)^j,
 \qquad k\geq0.
\]
This is the local analytic result of
Pham~\cite[Part~II, Section~1.1, (1.5)--(1.6) and the footnote on
printed page~26]{Pham}; it permits degenerate isolated critical
points. Pham credits Maisonobe and Rombaldi for the local solution
comparison used there. The amplitude $A e^{-q}$ is analytic. Every nonzero
term has a strictly negative power of $t$ and tends to zero after
removing $e^{\lambda t}$. The tails outside the local neighbourhood
are exponentially negligible by the uniform bound already proved.

The finite-critical decomposition accounts for the full module.
Group equal critical values before separating the distinct formal
exponential factors. Within each such factor the displayed local
expansions contain only decaying powers and logarithms. A nonzero
polynomial prefactor cannot occur. The zero function is allowed when
all these contributions vanish. This proves the theorem.
\end{proof}

\begin{corollary}[Original Gaussian markers]\label{cor:gradientmarkers}
For the standard Gaussian quadratic $q=\sum x_i^2/2$, suppose
\eqref{eq:gradientconditions} holds with $J=2$. Then every original
marked series $\Phi_A$ is either zero or is not a finite
exponential-polynomial sum. In particular $P$ is neither pure nor safe.
\end{corollary}
\begin{proof}
The exact top de Rham quotient and connection are
\[
 M=\C(t)[x]\Big/\sum_i(\partial_i+tP_i-x_i)\C(t)[x],
 \qquad\nabla[A]=[A_t+PA].
\]
Gaussian integration by parts makes $[A]\mapsto\Phi_A$ horizontal
on this vector-space quotient. The spatial and parameter derivatives
commute; no ordinary Jacobian-ideal quotient is substituted.
If some $\Phi_A$ is a
nonzero finite exponential-polynomial sum, a constant-coefficient
differential operator isolates one exponential and kills its
polynomial prefactor, giving $c e^{\beta t}$, $c\ne0$.
Because
\[
 \partial_t\Phi_A=\Phi_{PA},
\]
this is still an original fixed polynomial marker. Replace $P$ by
$P-\beta$ and divide the marker by $c$. Its formal seed value is one.
Lemma~\ref{lem:seedfunctional} gives an actual constant period for
that marker, contradicting Theorem~\ref{thm:gradientcriterion}.
Its comparison hypotheses hold with base $t$ and phase $tP-q$ on
$\mathbb A^1\times\mathbb A^n$, after removing the finite exceptional
base set. No convergence of the original Gaussian series is assumed.
Finally $\Phi_1(0)=1$, so it is a nonzero series; purity or safety
would make it a polynomial.
\end{proof}

\begin{corollary}[Smooth leading form, in every degree]
\label{cor:smoothleading}
If $\deg P=D\geq2$ and the homogeneous leading form satisfies
$\nabla F_D(v)=0$ only for $v=0$, then $P$ is not pure. In fact
Corollary~\ref{cor:gradientmarkers} holds for every original marker.
\end{corollary}
\begin{proof}
Compactness of the complex unit sphere gives
$\|\nabla P(x)\|\geq c\|x\|^{D-1}$ for large $\|x\|$.
Also $\|\operatorname{Hess}P\|=O(\|x\|^{D-2})$, so its ratio to
the squared gradient is $O(\|x\|^{-D})$. Both conditions for $J=2$
follow, with all lower terms retained.
\end{proof}

\begin{lemma}[The ordinary-node test]\label{lem:nodecoercivity}
Let $P=F_3+F_2+F_1+F_0$ in $n\geq2$ variables. Suppose its
projective singular set is finite and every singular point $[v]$
has $\ker\operatorname{Hess}F_3(v)=\C v$. If $F_2(v)\ne0$ at
every such point, then
$\|\nabla P(x)\|\geq c\|x\|$ outside a compact set. Consequently
$P$ is not pure under the Gaussian functional.
\end{lemma}
\begin{proof}
Otherwise choose an unbounded sequence with
$\|\nabla P(x)\|/\|x\|\longrightarrow0$. A projective limit $[v]$
is singular for $F_3$. Choose a fixed complement to $\C v$ and write
$x=sv+y$, where $y=o(s)$. With $H=\operatorname{Hess}F_3(v)$,
homogeneity gives the exact expansion
\[
 \nabla P(sv+y)=sHy+\nabla F_3(y)+s\nabla F_2(v)
                              +\nabla F_2(y)+\nabla F_1.
\]
The map $H$ is injective on the chosen complement. Divide by $s$.
The error from $\nabla F_3(y)/s$ is $o(\|y\|)$, so the equation
first forces $y$ to stay bounded. A limit $y_0$ satisfies
$Hy_0+\nabla F_2(v)=0$. Pair bilinearly with $v$. Symmetry and
$Hv=0$ kill the first term, whereas Euler's identity gives
$v\cdot\nabla F_2(v)=2F_2(v)$, a contradiction.
The Hessian of a cubic is $O(1+\|x\|)$, so the second condition
in \eqref{eq:gradientconditions} follows from the first. Apply
Corollary~\ref{cor:gradientmarkers}.
\end{proof}

\subsection{A common residual family for ordinary nodes}
\label{sec:ordinaryresidual}

For a pure cubic with only ordinary leading nodes, the preceding
lemma forces $F_2$ to vanish at some node. An isotropic such node
gives the frame already proved in Lemma~\ref{lem:frame}. Every
nonisotropic such node has the following single normal form.

\begin{lemma}\label{lem:ordinarynormal}
Let $[w]$ be an ordinary node of a ternary $F_3$, with $F_2(w)=0$
and $w^Tw\ne0$. An invertible affine change of coordinates puts
the pair consisting of $P$ and the Gaussian quadratic into
\begin{align*}
 P={}&T(uv+h)+a u^3+d v^3+b u^2+c v^2+e u+f v+k,\\
 q={}&\tfrac12(Au^2+Bv^2+mT^2+2Cuv+2DuT+2EvT)
                  +\ell_u u+\ell_v v+\ell_T T+q_0,
\end{align*}
where $m\ne0$ and $\det\operatorname{Hess}q\ne0$.
All Gaussian means and cross terms must be retained. Its leading
curve is irreducible nodal when $ad\ne0$, a transverse line and
smooth conic when exactly one of $a,d$ vanishes, and three
nonconcurrent lines when $a=d=0$.
\end{lemma}
\begin{proof}
Take $T$ in the direction $w$ and choose a complementary plane.
The singular-point equations make $F_3$ linear in $T$; its
coefficient is a nondegenerate binary quadratic because the node
is ordinary. Factor that quadratic and choose its two linear
factors as coordinates $u,v$. Thus $F_3=Tuv+B_3(u,v)$.
Replace $T$ by $T$ plus a linear form in $u,v$ to remove the
$u^2v$ and $uv^2$ terms of $B_3$.

There is no $T^2$ term in $F_2$, precisely because $F_2(w)=0$.
The full coefficient of $T$ is therefore
$uv+\alpha u+\beta v+h_0$. Translations of $u,v$ make it $uv+h$.
The cubic terms outside $T$ remain pure cubes. A constant shift
of $T$ removes the remaining quadratic $uv$ term. This gives the
displayed $P$. The direction of $T$ has not changed, so $m\ne0$;
invertibility of the affine change preserves the nondegenerate
Gaussian Hessian, while translations and the shifts of $T$ produce
the means and cross terms shown.

If $ad\ne0$, the polynomial $Tuv+a u^3+d v^3$ is irreducible:
an independent factor of a polynomial linear in $T$ would divide
both $uv$ and $a u^3+d v^3$, whose gcd is one. Its gradient shows
that its sole singular point is the stated ordinary node. If
$a=0,d\ne0$, it is $v(Tu+d v^2)$, with a smooth conic meeting
the line transversely at two points; the other one-zero case is
symmetric. If both vanish, it is the coordinate triangle $Tuv$.
\end{proof}

This consolidates three leading-curve rows into one residual family.
It is a reduction, not yet an exclusion of that family. In particular
the full quadratic $q$, including its means, cannot be replaced by
a convenient diagonal form during the remaining period argument.

\subsection{Two controls on the scope of the criterion}

Ordinary tameness of $P$ alone is insufficient for the Gaussian
conclusion. Let $P=\ell^Tx$, with $\ell=(1,i)/\sqrt2$ in two
variables. Its gradient is constant and nonzero, but
$\Phi_1=\exp(t^2\ell^T\ell/2)=1$. The phase $tP-q$ has its
critical point at $x=t\ell$, escaping with action zero. Thus an
isolated critical set for $P$, or the weaker gradient bound away
from zero, would lose exactly the contribution that purity needs.
Condition~\eqref{eq:gradientconditions} with $J=2$ prevents this.

The distinction between an original marker and a time-dependent
amplitude also remains essential. For
$P=(X_1^2+X_2^2)/2$ one has $\Phi_1=(1-t)^{-1}$ and
$(1-t)\Phi_1=1$. This phase satisfies the gradient criterion.
The unit in its rational differential module is represented by a
time-dependent amplitude. No original fixed polynomial marker has
a nonzero polynomial period, by Corollary~\ref{cor:gradientmarkers}.
The abstract module alone does not record that distinction.

\subsection{Ordinary nodes with a finite affine critical set}
\label{sec:ordinaryfinite}

The three ordinary-node families admit one argument whenever the
critical set of $P$ is finite. In particular we need not solve for
each escaping critical point, assume it is Morse, or separate its
action from the actions of other escaping points. The Gaussian
quadratic itself detects the unwanted contribution at infinity.

Retain the notation and the full pair $(P,q)$ of
Lemma~\ref{lem:ordinarynormal}. Write $q_2$ for the homogeneous
quadratic part of $q$. The absence of a Gaussian frame is exactly
the following condition at the leading nodes:
\begin{equation}\label{eq:nodenoframe}
 F_2(w)=0\quad\Longrightarrow\quad q_2(w)\ne0.
\end{equation}
Indeed the Gaussian bilinear form in these affine coordinates is
$2q_2$; the value of $F_2$ at a leading singular direction is
unchanged by translation. In the common normal form the $T$ node
already satisfies this condition because $m\ne0$. The only extra
conditions are $c\ne0$ or $B\ne0$ when $d=0$, and $b\ne0$ or
$A\ne0$ when $a=0$.

\begin{lemma}[The stationary action identity]\label{lem:actionidentity}
Put $r=1/t$, and let $x(r)$ be an algebraic critical branch of
$S=P/r-q$. A prime means total differentiation along the branch.
Then
\begin{equation}\label{eq:actionidentity}
 S'=-P(x(r))/r^2,\qquad q(x(r))=-S-rS'.
\end{equation}
If $x(r)$ is bounded, the polar part of $S$ is $\alpha/r$, where
$\alpha=P(x(0))$. If it escapes and \eqref{eq:nodenoframe} holds,
its polar part is neither zero nor $\alpha/r$ for any constant
$\alpha$.
\end{lemma}
\begin{proof}
Use a finite Puiseux cover so every branch is meromorphic. At a
critical point the chain-rule term $\nabla_x S\cdot x'$ vanishes,
giving both identities. For a monomial $c r^{-\kappa}$ of the
polar part of $S$, its contribution to $-S-rS'$ is
$(\kappa-1)c r^{-\kappa}$. Thus bounded $q(x(r))$ permits only
the pole $\alpha/r$ in $S$. Its coefficient is $P(x(0))$ when
$x$ is bounded.

For an escaping branch let $R=\|x\|$. The critical equation is
$\nabla P=r\nabla q$, so $\|\nabla P\|/R\longrightarrow0$.
The leading projective direction is a singular point of $F_3$.
The ordinary-node argument in the proof of the preceding coercive
lemma works at this one direction: write $x=sw+y$, $y=o(s)$,
in a fixed complementary plane. The injectivity there of
$H=\operatorname{Hess}F_3(w)$ forces $y$ bounded from
\[
 \nabla P/s=Hy+\nabla F_2(w)+O((1+\|y\|)^2/|s|)=o(1).
\]
Pairing with $w$ then gives $F_2(w)=0$. Condition
\eqref{eq:nodenoframe} gives $q_2(w)\ne0$, so $q(x(r))$ has a
pole: its leading quadratic term cannot be cancelled by its
lower-degree terms. Identity~\eqref{eq:actionidentity} now rules
out both claimed polar parts. This argument includes all escape
scales and all coincident actions.
\end{proof}

We specify the classical rank input before using it. This is a
convergent algebraic application of formal stationary phase, not
an extension to arbitrary nonconvergent coefficients or rings.

\begin{lemma}[Rank accounting for a tame algebraic pencil]
\label{lem:tamepencilrank}
For the present no-frame pencil, shrink a punctured $r$ disk to
avoid its finitely many exceptional parameters. Its top twisted
de Rham module has rank
\[
 N=\sum_{\text{critical points of }S_r}\mu.
\]
After a finite cover, the rank of its block with nonzero
exponential polar part $\gamma$ is the sum of these generic local
Milnor numbers over critical branches with action polar part
$\gamma$. Equal actions add their multiplicities. Consequently
the rank of the zero-exponential block is the sum of the generic
Milnor numbers of the branches with bounded action.
\end{lemma}
\begin{proof}
For fixed small nonzero $r$, the homogeneous quadratic term of
$S_r$ at a leading node $w$ is $F_2(w)/r-q_2(w)$. It is nonzero
at every node, by \eqref{eq:nodenoframe} and exclusion of finitely
many parameters. The ordinary-node coercivity lemma gives
$\|\nabla S_r\|\ge c_r\|x\|$ outside a compact set. The
critical locus is therefore finite: a positive-dimensional affine
algebraic component would be unbounded. Pham's tame Milnor
decomposition~\cite[Part~I, Sections~1--4]{Pham} gives the displayed
rank, retaining the Milnor number of every degenerate point.

Here are the checks needed to apply the direct-image rank formula.
Take $U=\C_r^*\times\C_x^3$, $f=r$, $g=S_r$, and the regular
algebraic module $\mathcal O_U$. Form the ordinary direct image
by $(f,g)$, followed by exponential integration in the $g$
coordinate. All these are algebraic maps; $g$ is regular on $U$.
The ordinary direct image is regular holonomic. Over the punctured
$r$ disk its relevant cohomology is the middle, reduced
Gauss--Manin module of the fibers of $S_r$. The other, constant
fiber cohomology has zero exponential integral in the affine $g$
line. Punctual parameter terms disappear on the punctured disk.

There are no extra finite-$g$ singular curves from spatial infinity.
The coercive estimate is locally uniform on a compact subset of
the punctured $r$ disk. For $g$ in a compact set, the identity
$P=r(g+q)$ on a fiber gives
\[
 \partial_r S_r=-(g+q)/r=O(1+\|x\|^2).
\]
Lifting $r$ variation by the normalized gradient therefore uses a
vector field of at most linear growth. Lifting $g$ variation uses
a field $O(1/\|x\|)$. These fields are complete outside a fixed
ball on each such parameter neighborhood. They trivialize the
outer fiber pairs, so all nonvertical finite-$g$ discriminant
curves come from affine critical values. Along such a curve the
conormal multiplicity of the middle Gauss--Manin module is the
sum of the local Milnor numbers: a transverse target disk computes
precisely the corresponding isolated vanishing cycles. This also
applies to persistent degenerate points and equal critical values.

First make a common finite Puiseux base change of all the
critical-value curves. Apply the nonramified case of Roucairol's
Theorem~0.1(2), proved in her Section~1.2, and Theorem~0.3
\cite{RoucairolFormal} to this identified ordinary direct image.
They give exactly the stated nonzero polar parts and ranks.
Count every graph after this base change, including distinct
conjugates with the same polar part. For example the two graphs
$g=r^{-1}\pm r^{1/2}$ must contribute twice after $r=\rho^2$;
retaining just their one polar expression would lose rank.
Their rank statement requires no separation of the constant terms
of equal polar parts. The additional separation hypothesis in her
Theorem~0.2 concerns its monodromy formula, which we do not use.
Vertical components at $r=0$ and the boundary $g=\infty$ do not
add polar branches in Theorem~0.1. Subtracting the nonzero-block
ranks from $N$ proves the last assertion. In particular we have
not inferred the regular rank from an Euler characteristic alone.
\end{proof}

\begin{proposition}\label{prop:ordinaryfinite}
Suppose the common ordinary-node pair satisfies
\eqref{eq:nodenoframe} and $P$ has finite affine critical set.
Then $P$ is not pure. More generally every original Gaussian
marked period is either zero or is not a finite sum of
polynomials in $t$ times exponentials $e^{\beta t}$.
\end{proposition}
\begin{proof}
Every bounded critical branch of $S_r$ limits to a critical point
of $P$. Disjoint small Milnor neighborhoods of these finitely many
points contain, for the deformation $P-rq$, total critical
multiplicity $\mu_p$ in the neighborhood of $p$. This is
conservation of the local intersection multiplicity of the
gradient, and does not require a Morse perturbation. By
Lemma~\ref{lem:actionidentity} its action has polar part
$P(p)/r$. Every escaping branch has nonzero polar part. Hence
Lemma~\ref{lem:tamepencilrank} identifies the entire zero-exponential
rank as
\begin{equation}\label{eq:ordinaryzerorank}
 \sum_{p:\,dP(p)=0,\,P(p)=0}\mu_p.
\end{equation}

We construct that many independent local formal functionals and
compute their monodromy over the unramified field $\C((r))$.
At an isolated critical point at least one of $u,v$ is nonzero.
Indeed if both vanish, criticality requires $h=e=f=0$, in which
case the whole $T$ axis is critical. In the chart $u\ne0$ put
\[
 s=uv+h,\quad R(u)=G(u,-h/u),\quad
 T'=T+\frac{G(u,(s-h)/u)-G(u,-h/u)}s,
\]
where $G=P-T(uv+h)$. The quotient is analytic, and exactly
$P=sT'+R(u)$. The coordinate Jacobian is $u\ne0$. At a critical
point $p$ the one-variable critical multiplicity of $R$ is
$\mu_p$. An analytic one-variable coordinate $y$, with
$w=T'$, therefore gives the exact identity
\[
 P=\alpha+sw+y^k,\qquad k=\mu_p+1,\quad\alpha=P(p).
\]
The volume Jacobian and $e^{-q}$ multiply every amplitude by an
analytic unit $V(y,s,w)$; no Gaussian mean or cross term has been
discarded.

For $j=1,\ldots,k-1$ define a formal local functional on an
analytic monomial $y^a s^b w^c$ as zero unless $b=c$ and
$a=j-1+\ell k$ for an integer $\ell\ge0$. Otherwise its value is
\begin{equation}\label{eq:ordinarycharacters}
 e^{\alpha/r}r^{1+j/k}
       (-1)^{b+\ell}b!\,(j/k)_\ell\,r^{b+\ell}.
\end{equation}
Evaluate a polynomial form after multiplying its transformed
amplitude by $V$. At each power of $r$ only finitely many Taylor
terms contribute. The recurrences
\[
 I_{a+k}=-\frac{a+1}{k}r I_a,\qquad
 J_{b+1,b+1}=-(b+1)rJ_{b,b}
\]
prove spatial integration by parts and parameter horizontality.
Thus these are functionals on the full twisted quotient, not on
an assumed cyclic image. They exist even if the corresponding
asymptotic series diverge; their input coefficients here are the
Taylor series of the stated analytic coordinates.

They are independent. Prescribe, at each $p$, the jets
$1,y,\ldots,y^{k-2}$ on the line $s=w=0$ and zero jets at the
other points, using polynomial representatives of finite local
jets and the Chinese remainder theorem. After dividing column
$j$ by $e^{\alpha/r}r^{1+j/k}$, the leading evaluation matrix is
the jet map multiplied by the analytic unit $V$. It is triangular
with nonzero diagonal. Its higher terms are in $r\C[[r]]$.
For $\alpha=0$ these independent characters fill exactly the rank
in \eqref{eq:ordinaryzerorank}; there is no unaccounted regular
piece. Their formal monodromy eigenvalues, or their inverses if
one uses the dual module, are
\[
 \exp(2\pi i j/k),\qquad 1\le j<k.
\]
None is one. This conclusion is over $\C((r))$. We do not replace
a noninteger residue by its integer multiple after ramification.
The unique zero-exponential projector commutes with descent, and
the displayed unramified coefficient matrices identify its whole
block. Every other block has a nonzero exponential polar part.
The ambient module therefore has no trivial composition factor.

For completeness this excludes the formal Gaussian unit without
choosing an actual analytic period. Formal Gaussian integration
annihilates $\partial_i+tP_i-q_i$ coefficientwise. If its value
on a fixed marker $A$ were a nonzero constant, restriction to the
rational differential submodule generated by $A$ would give a
surjection to the trivial module: apply the functional to every
rational relation between $A,\nabla_tA,\ldots$. A relation has
zero constant coefficient because all positive derivatives of
that constant vanish. Exact formal extension at infinity would
preserve this composition factor, contradicting what was proved.

Finally a nonzero finite exponential-polynomial can be reduced
to $c e^{\beta t}$ by a constant-coefficient differential
operator. On an original marker this operator replaces $A$ by
the fixed polynomial marker $T(P)A$. Replace $P$ by $P-\beta$
and divide the marker by $c$. The hypotheses and finite critical
set are unchanged, so the same contradiction applies. Taking
$A=1$ proves the purity exclusion.
\end{proof}

\subsection{An exhaustive critical-set split}

The finite-critical hypothesis in Proposition~\ref{prop:ordinaryfinite}
has a short exhaustive coefficient test. This prevents a generic
argument from silently discarding the positive-dimensional cases.

If $h\ne0$, the critical locus lies on $uv=-h$, and the equations
above identify it with $R'(u)=0$ in $\C^*$. The Laurent polynomial
\[
 R(u)=a u^3+b u^2+e u+k-fh/u+ch^2/u^2-dh^3/u^3
\]
is nonconstant unless $a=b=c=d=e=f=0$. A nonzero Laurent
derivative has only finitely many roots in $\C^*$. In the constant
case the critical set is the smooth conic $uv=-h$, $T=0$.

If $h=0$, outside the center the critical equations are
\[
 \begin{array}{lll}
 v=0,\ u\ne0:&3au^2+2bu+e=0,&T=-f/u,\\
 u=0,\ v\ne0:&3dv^2+2cv+f=0,&T=-e/v.
 \end{array}
\]
The center $u=v=0$ is critical exactly when $e=f=0$, and then
the whole $T$ axis is critical. Consequently the following table
is exhaustive and disjoint:
\begin{center}\small
\begin{tabular}{@{}p{0.51\textwidth}p{0.39\textwidth}@{}}
\toprule
Condition & Conclusion\\
\midrule
$h\ne0$, $(a,b,c,d,e,f)\ne0$ & Finite critical set; excluded above\\
$h\ne0$, $a=b=c=d=e=f=0$ & Critical conic; Corollary~\ref{cor:ordinaryconics}\\
$h=0$, $e=f=0$ & Critical center axis, possibly extra curves; open here\\
$h=0$, $(e,f)\ne0$, $a=b=e=0$ & Critical hyperbola; Corollary~\ref{cor:ordinaryconics}\\
$h=0$, $(e,f)\ne0$, $d=c=f=0$ & Symmetric hyperbola; same corollary\\
$h=0$, $(e,f)\ne0$, $(a,b,e)\ne0$, $(d,c,f)\ne0$
 & Finite critical set; excluded above\\
\bottomrule
\end{tabular}
\end{center}
The two hyperbola rows cannot overlap because $(e,f)\ne0$.
All isolated degeneracies in the finite rows have already been
included through their Milnor multiplicities. The conic and hyperbola rows are proved in the next subsection.
The center-axis row still needs a full regular-block or homology
argument; the action identity alone does not settle it.

\subsection{A critical conic in every polynomial degree}
\label{sec:hyperbolicsuspension}

The conic and hyperbola rows of the preceding table are manifestations
of the same calculation. It also works when one of the cubic terms is
replaced by a polynomial of arbitrary degree. We prove that stronger
statement to make clear what the mechanism actually uses.

\begin{theorem}\label{thm:hyperbolicsuspension}
Retain a full nondegenerate affine Gaussian quadratic
\[
 q=\tfrac12(Ax^2+By^2+mw^2+2Cxy+2Dxw+2Eyw)
                     +\ell_x x+\ell_y y+\ell_w w+q_0.
\]
Let $AB\ne0$ and
\[
 P=wxy+F(w)+k,\qquad F(0)=0,\quad F'(0)=h\ne0.
\]
Allow any polynomial degree $d_F=\deg F\ge1$, requiring $m\ne0$
when $d_F=1$. Then $P$ is not pure. Every original marked Gaussian
period is either zero or is not a finite exponential-polynomial.
The generic top twisted rank is $d_F+3$ for $d_F\ge2$ and five for
$d_F=1$. No Morse assumption on the Laurent phase obtained by
restricting $q$ to the conic, or on the isolated critical points of
$P$, is needed.
\end{theorem}

The nonzero quadratic poles at both ends of the conic are essential
to the rank-four construction below. The condition $AB\ne0$ states
this explicitly. Nondegeneracy of the full Gaussian alone would not
imply it. This theorem concerns complex polynomials and valued-field
formal frames; it does not assert an all-ring extension.

\begin{proof}
We first justify the full rank. Fix a generic $t\ne0$ and set
$S=tP-q$. Its gradient has a linear lower bound at complex spatial
infinity. To see this, suppose along an unbounded sequence
$\|\nabla S\|=o(R)$, where $R=\|(x,y,w)\|$. If $w$ stays
bounded, pass to a limit $w_0$ and normalize $(x,y)$ by $R$.
The first two gradient equations give
\[
 -AX+(tw_0-C)Y=0,\qquad (tw_0-C)X-BY=0.
\]
At least one of $X,Y$ is nonzero, and $AB\ne0$ makes both
nonzero. But the last gradient equation, divided by $R^2$, gives
$tXY=0$, a contradiction.

If $w$ is unbounded, the first two equations have a matrix with
inverse $O(1/|w|)$. Hence
$x,y=O(1)+o(R)/|w|$. First $R\asymp |w|$, then $x,y=O(1)$.
The last equation is therefore
\[
 tF'(w)-mw+O(1)=o(|w|).
\]
This is impossible for $d_F\ge3$, for $d_F=2$ with
$2t[w^2]F-m\ne0$, and for $d_F=1$ with $m\ne0$.
The same proof is uniform on compact subsets of these coefficient
conditions. Thus all phase-critical points stay in one ball during
any compact coefficient deformation inside that domain. They are
isolated, and their total Milnor number is constant by conservation
of local gradient-intersection multiplicity. Pham's tame Milnor
decomposition~\cite[Part~I, Sections~1--4]{Pham} identifies that
number with the full twisted rank.

For fixed $t$, the coefficient domain just used is the complement
of finitely many complex algebraic hypersurfaces, and is path
connected. We may compute the number at a generic diagonal
Gaussian with zero means, retaining $ABm\ne0$, and a generic $F$.
The critical equations there are
\[
 twy=Ax,\qquad twx=By,\qquad t(xy+F'(w))=mw.
\]
There are $d_F-1$ simple solutions with $x=y=0$ if $d_F\ge2$,
and one if $d_F=1$. The other solutions have
$w=\pm\sqrt{AB}/t$. At each of these two values, $x,y$ are
proportional and their nonzero product is prescribed by the last
equation; generically this gives two simple solutions. The counts
are $d_F+3$ and five, respectively. Uniform coercivity and the
connected deformation justify this specialization; a convenient
critical-point count by itself would not justify the rank for a
degenerate member.

We now construct a full formal frame. Set $r=t^{-1/2}$ and write
$F(w)=hw+bw^2+\sum_{j\ge3}f_jw^j$. Near the entire critical
conic $w=0$, $xy=-h$, use the exact coordinates
\[
 y=(rV-h)/x,\qquad w=rU.
\]
They give
\begin{align}
 tP-q&=k/r^2+UV+bU^2+\psi(x)+E_r(x,V,U),
                  &E_r\in r\C[x,x^{-1},V,U][r],\label{eq:conicphase}\\
 dx\,dy\,dw&=r^2\,\frac{dx}{x}\,dV\,dU,\label{eq:conicvolume}\\
 \psi(x)&=-\frac A2x^2-\frac{Bh^2}{2x^2}
                -\ell_x x+\frac{h\ell_y}{x}+Ch-q_0.\notag
\end{align}
All Gaussian cross terms and means are in $\psi$ or $E_r$.
The normal quadratic $UV+bU^2$ has Hessian determinant $-1$ for
every $b$. Its normalized formal Gaussian functional is well
defined and takes one to one.

The logarithmic Laurent quotient for $\psi$ has rank four, with
basis $1,x,x^2,x^3$. Here is a direct proof, including independence.
Its relation on $x^j$ is
\[
 jx^j-Ax^{j+2}-\ell_xx^{j+1}
             -h\ell_yx^{j-1}+Bh^2x^{j-2}=0.
\]
The top coefficient eliminates all exponents at least four without
creating a negative exponent; the bottom eliminates all negative
exponents without creating an exponent above three. If the image
of a nonzero Laurent polynomial under $x\partial_x+x\psi_x$
were supported in $[0,3]$, its highest exponent would be at most
one and its lowest at least two, which is impossible.

Let $\lambda_0,\ldots,\lambda_3$ be the dual basis on this
quotient, tensored with the normal Gaussian functional. For a
polynomial amplitude $a(x,y,w)$, define
\begin{equation}\label{eq:coniccharacters}
 L_i(a)=e^{k/r^2}r^2\lambda_i\!\left(
    e^{E_r}a(x,(rV-h)/x,rU)\right).
\end{equation}
Every coefficient is a finite Laurent-polynomial calculation, so
this is a well-defined formal series. Multiplication by $e^{E_r}$
conjugates the fixed twisted differential to the perturbed one.
Consequently these four functionals annihilate all spatial exact
forms. The moving volume is retained in the factor $r^2$.
For an explicit parameter check the coordinate velocity in the
original variables is
\[
 Z=\left(0,\frac{xy+h}{rx},\frac wr\right),
                   \qquad \operatorname{div}Z=2/r.
\]
Cartan's identity $\mathcal L_Z=d\iota_Z+\iota_Zd$ shows that
the pullback derivative differs from the original parameter
connection only by a spatial exact form and this volume
derivative. Thus \eqref{eq:coniccharacters} are horizontal for
the original full module. In particular, for every fixed
polynomial $a$,
\begin{equation}\label{eq:conicseedorder}
 L_i(a)\in e^{kt}r^2\C[[r]].
\end{equation}
This bound holds for all original markers, not just for the unit.

Suppose first $d_F\ge2$. The other critical points of $P$ are
$x=y=0$, $F'(w_p)=0$. Each $w_p$ is nonzero because $h\ne0$.
If it has multiplicity $\mu_p$, analytic coordinates give
\[
 P=k+F(w_p)+s z+y_0^{\mu_p+1}.
\]
Indeed absorb the nonvanishing factor $w$ into $x$ or $y$, then
use the one-variable power coordinate for $F$ at $w_p$.
The explicit characters \eqref{eq:ordinarycharacters}, including
their integration-by-parts and time recurrences, apply to this
exact local model. Their prefactors are
\[
 e^{(k+F(w_p))t}t^{-1-j/(\mu_p+1)},
                    \qquad 1\le j\le\mu_p.
\]
The Gaussian and coordinate Jacobian are analytic units in that
construction. These characters retain multiple roots and equal
critical values. Their total number is $\sum_p\mu_p=d_F-1$.

The $d_F+3$ characters just written are independent on polynomial
amplitudes. Put $B_0(w)=\prod_p(w-w_p)^{\mu_p}$. For the four
conic rows use $x^jB_0(w)$, $0\le j\le3$. Their leading
conic matrix is $B_0(0)I_4$. Their leading finite jets vanish.
For the remaining rows choose polynomials in $w$ with value zero
at $w=0$ and the separate Hermite jets of orders
$0,\ldots,\mu_p-1$ at the nonzero $w_p$. The Chinese remainder
theorem supplies them. Their conic values have an additional
positive power of $r$. Their normalized finite evaluation matrix
is the invertible local jet map, multiplied by analytic units.
After dividing each character column by its scalar prefactor,
the leading matrix is therefore block diagonal and invertible.
Since its size is the proved full rank, these are a full formal
frame. No uncomputed image kernel or infinity block remains.

For $d_F=1$ there is instead one further character. In the exact
coordinates
\[
 x=X/t,\qquad y=Y/t,\qquad w=(h/m)t+W,
\]
the phase is
\begin{align*}
 tP-q={}&\frac{h^2}{2m}t^2+(k-h\ell_w/m)t+H(X,Y,W)+O(t^{-1}),\\
 H={}&\frac hm XY-\frac hm(DX+EY)-\frac m2 W^2-\ell_w W-q_0.
\end{align*}
The correction is polynomial in $t^{-1},X,Y,W$, with positive
$t^{-1}$ order. The Hessian determinant of $H$ is $h^2/m\ne0$.
Formal Gaussian evaluation gives a horizontal character with
nonzero quadratic action and volume prefactor $t^{-2}$.
To check independence set $w_*=(h/m)t$. Use
$x^j(1-w/w_*)$, $0\le j\le3$, and $w/w_*$ as five polynomial
amplitudes with rational parameter coefficients. The former
have the original conic leading matrix and zero leading escape
values; the latter has zero leading conic values and nonzero
leading escape value. This gives the full rank-five frame.

In either degree case, after finite ramification the frame splits
the module into rank-one factors. Every zero-exponential
functional has positive order on every original polynomial
marker, by \eqref{eq:conicseedorder} and the displayed finite
prefactors. Lemma~\ref{lem:positiveformal} rules out a nonzero
constant original marked period. A constant-coefficient operator
reduces any nonzero finite exponential-polynomial to $c e^{\beta t}$;
it replaces the marker by another original polynomial in $P$.
Shifting $k$ by $-\beta$ and normalizing the marker gives the same
contradiction. This proves the theorem.
\end{proof}

\begin{corollary}\label{cor:ordinaryconics}
All the critical-conic and critical-hyperbola rows in the ordinary-node
table are not pure. Only its center-axis row remains open here.
\end{corollary}
\begin{proof}
For $h\ne0$ and $a=b=c=d=e=f=0$ the residual is
$T(uv+h)+k$. Apply the $d_F=1$ case with $w=T$, $x=u$, $y=v$.
The no-frame conditions give $ABm\ne0$.

For $h=0$, $a=b=e=0$, $f\ne0$, permute the coordinates and
write the residual as
\[
 v(Tu+d v^2+c v+f)+k.
\]
Take $w=v$, $x=T$, $y=u$, and
$F(w)=d w^3+c w^2+f w$. The quadratic coefficients along $T,u$
are nonzero: the original $T$ square coefficient is nonzero, and the $u$
node has $F_2=0$, so its Gaussian square coefficient is nonzero
by the no-frame condition. If this $F$ is linear, the remaining
$v$ node also has $F_2=0$, giving the extra $w$ square coefficient
required in degree one. The symmetric hyperbola row is identical
after swapping $u,v$. These are all the indicated rows.
\end{proof}


\subsection{The geometric case split for frame extraction}

For a cubic of degree exactly three, either the frame test
\eqref{eq:frame} succeeds, or it does not. In the latter case,
the projective curve $F_3=0$ belongs to exactly one of these types:
\begin{center}\small
\begin{tabular}{@{}p{0.33\textwidth}p{0.56\textwidth}@{}}
\toprule
Leading curve & Proof or remaining no-frame exclusion\\
\midrule
Smooth cubic & Proved by Corollary~\ref{cor:smoothleading}; no pure phase\\
Irreducible nodal cubic & Finite-critical no-frame members excluded by Proposition~\ref{prop:ordinaryfinite}; center-axis cases remain\\
Irreducible cuspidal cubic & Smooth parabola, separated lines and double-line critical profiles\\
Line and smooth conic, transverse & Same finite-critical exclusion; center-axis cases remain\\
Line and smooth conic, tangent & Tangency, degenerate finite blocks and escaping profiles\\
Three nonconcurrent lines & Same finite-critical exclusion; center-axis cases remain\\
Three concurrent distinct lines & Constant and nonconstant transverse linear coefficient\\
Double line and a distinct line & Isotropic and nonisotropic repeated-line directions\\
Triple line & Isotropic and nonisotropic repeated-line directions\\
\bottomrule
\end{tabular}
\end{center}
The list is exhaustive by factorization of a plane cubic over $\C$.
The smooth row is now proved here. The ordinary-node gate and its
common residual normal form are also proved. Proposition~\ref{prop:ordinaryfinite}
excludes every finite-critical no-frame member of that family; the
conic and hyperbola cases are excluded by Corollary~\ref{cor:ordinaryconics};
only the center-axis cases remain open in this family. The other rows remain
proof obligations. For isolated leading singular points, a null
point with $F_2=0$ goes directly to the affine-linear-coordinate table;
otherwise the lower Gaussian quadratic and its cross terms must be
retained in the indicated exclusion. In the repeated-line cases the
singular locus is a whole projective line, and the frame test must be
applied to all its null directions, not to an arbitrarily selected point.

Thus the routing is now explicit. Completing frame extraction still
requires the ordinary-node center-axis exclusions and every other open
row, including a full rank or homology-completeness argument. The full affine-linear proof above does not prove those geometric
exclusions. Finite algebra checks do not replace their missing
rank and homology arguments.


\subsection{How the two reductions finish the theorem}
Proposition~\ref{prop:wlinear} is proved above. Assume the remaining
frame-extraction statement, Proposition~\ref{prop:frameextract}.
For $\deg P\leq2$, Proposition~\ref{prop:quadratic} gives the result.
For degree three, frame extraction and Lemma~\ref{lem:frame} put $P$
in the family of Proposition~\ref{prop:wlinear}. A complex orthogonal
linear substitution preserves both total degree and Gaussian moments.
The bound therefore transfers back to every original multiplier $Q$.
Equation~\eqref{eq:sharp} proves optimality. This is the entire logical
route from the two reductions to Theorem~\ref{thm:cubic}; a finite
pure-moment prefix is not needed.

\section{Completing the boundary deduction}
\begin{proof}[Proof of Theorem~\ref{thm:boundary}]
Wilson proves the assertion for $n\leq2$ in every degree.
Proposition~\ref{prop:quadratic} proves it for $d\leq2$ in every
dimension. Theorem~\ref{thm:cubic} supplies $(n,d)=(3,3)$.

Conversely, suppose none of these conditions holds. Then $n\geq3$
and $d\geq3$. If $n=3$, the excluded pair $(3,3)$ leaves $d\geq4$,
and $P_3$ from Section~\ref{sec:failures} is a counterexample.
If $n\geq4$, use $P_4$. Additional independent Gaussian variables can
be ignored, and increasing the allowed degree does not remove an
existing example. These possibilities exhaust the complement.
\end{proof}

The classification therefore has two minimal failures: degree four
in three variables, and degree three in four variables. The ternary
cubic theorem identifies the last positive corner between them.
The explicit threshold in that corner tells us more than eventual
vanishing: a multiplier of degree $q$ can never expose a nonzero
moment beyond power $3q$, and the shear shows that power $3q$ itself
can still occur.

\section*{What remains before submission}
The dimension--degree deduction, quadratic argument, all-order
counterexamples and sharpness proof are complete within their stated
inputs. The cubic algebraic subcase, Airy reduction, affine $WT$ face,
formal seed lemma and full mixed-chart formal proof are now written out.
The general stationary-ball construction gives a full formal frame
for every one-variable polynomial phase. Its leading-coefficient
criterion closes the entire rank-one $WT^2$ family, including the
flattened quadratic boundary, without an analytic shrinking estimate.
The affine-linear case table and its proof are now complete.
A marked semistable model, explicit nearby-cycle hypotheses and
logarithmic annulus and tail bounds supply the full Laurent estimate,
including logarithmic extensions and regular punctures. This proves
Proposition~\ref{prop:wlinear} within the manuscript. Frame extraction
is now proved for smooth leading forms, and its ordinary-node rows
share one explicit residual normal form after the proved coercive gate.
The finite-critical part of that residual is now excluded, including all
its isolated degeneracies and coincident actions. The action identity
prevents any escaping branch from hiding in the regular block; the
complete regular rank is accounted for using the precisely credited
nonramified direct-image rank theorem after a finite base change.
A common all-degree theorem now closes the critical conic and
hyperbola cases, with the full rank proved by uniform coercivity
and coefficient deformation. Only center axes remain in the
ordinary-node family; the other five geometric families also remain.

The perfect period comparison is used for representation and
completeness, not as a quantitative bound. The latter is proved
separately on compact pieces, necks and descending tails. Finite
symbolic identities do not establish these analytic statements or
the remaining geometric no-frame exclusions. The research archive's
internal theorem status is unchanged, but neither
submission readiness nor an independent mathematical review is claimed.
Completing these proofs remains the active paper objective.

\clearpage
\begin{thebibliography}{99}
\bibitem{DVEZ}
H.~Derksen, A.~van den Essen and W.~Zhao,
\emph{The Gaussian Moments Conjecture and the Jacobian Conjecture},
\emph{Israel Journal of Mathematics} \textbf{219} (2017), no.~2,
917--928; \href{https://arxiv.org/abs/1506.05192}{arXiv:1506.05192}.

\bibitem{DLMF}
NIST Digital Library of Mathematical Functions,
\href{https://dlmf.nist.gov/10.25.E1}{10.25.1} and
\href{https://dlmf.nist.gov/10.31.E2}{10.31.2},
modified Bessel equation and logarithmic power series.

\bibitem{HienRoucairol}
M.~Hien and C.~Roucairol, \emph{Integral representations for solutions
of exponential Gau\ss--Manin systems},
\emph{Bulletin de la Soci\'et\'e Math\'ematique de France}
\textbf{136} (2008), no.~4, 505--532,
\href{https://doi.org/10.24033/bsmf.2564}{doi:10.24033/bsmf.2564}.

\bibitem{Liu}
Z.~Liu, \emph{Sharp finite Gaussian moments for quadratic polynomials},
version~3, public preprint (2026),
\href{https://doi.org/10.5281/zenodo.21688251}{doi:10.5281/zenodo.21688251}.

\bibitem{Long}
C.~D.~Long, \emph{Small Counterexamples to the Gaussian Moments Conjecture},
public preprint (2026),
\href{https://arxiv.org/abs/2607.18186v1}{arXiv:2607.18186v1}.

\bibitem{Pham}
F.~Pham, \emph{La descente des cols par les onglets de Lefschetz,
avec vues sur Gauss--Manin}, \emph{Ast\'erisque} \textbf{130}
(1985), 11--47,
\href{https://numdam.org/item/AST_1985__130__11_0/}{Numdam primary text}.

\bibitem{RoucairolFormal}
C.~Roucairol, \emph{Formal structure of direct image of holonomic
$\mathcal D$-modules of exponential type},
\emph{Manuscripta Mathematica} \textbf{124} (2007), no.~3, 299--318;
\href{https://arxiv.org/abs/math/0604134v2}{arXiv:math/0604134v2}.

\bibitem{SabbahStationary}
C.~Sabbah, \emph{An explicit stationary phase formula for the local
formal Fourier--Laplace transform}, in \emph{Singularities I},
Contemporary Mathematics \textbf{474}, American Mathematical Society,
2008, 309--330;
\href{https://arxiv.org/abs/0706.3570v2}{arXiv:0706.3570v2}.

\bibitem{Wall}
C.~T.~C.~Wall, \emph{Singular Points of Plane Curves}, Cambridge
University Press, 2004, Chapter~2, 15--38,
\href{https://doi.org/10.1017/CBO9780511617560.003}{doi:10.1017/CBO9780511617560.003}.

\bibitem{Wilson}
M.~Wilson, \emph{A face-isolation proof of the two-variable Gaussian
Moments Conjecture}, public preprint (2026),
\href{https://arxiv.org/abs/2607.23887v1}{arXiv:2607.23887v1}.
\end{thebibliography}
\end{document}
